\documentclass[11pt,a4paper]{article}
\usepackage[utf8]{inputenc}
\usepackage[english]{babel} % Supports both German and English sections
\usepackage{amsmath}
\usepackage{amssymb}
\usepackage{amsfonts}
\usepackage{amsthm}
\usepackage{graphicx}
\usepackage{booktabs}
\usepackage{geometry}
\usepackage{microtype}
\usepackage{array}
\usepackage{hyperref}

\geometry{margin=1in}

\hypersetup{
	colorlinks=true,
	linkcolor=blue,
	filecolor=magenta,      
	urlcolor=cyan,
	pdftitle={The E8 Lattice: Radial Topology, Zeta Holonomy, and the Algebraic Symmetrieleiter - Version 7},
	pdfpagemode=FullScreen,
}

\newtheorem{theorem}{Theorem}[section]
\newtheorem{lemma}[theorem]{Lemma}
\theoremstyle{definition}
\newtheorem{definition}[theorem]{Definition}
\newtheorem{remark}[theorem]{Remark}

\title{\textbf{The $E_8$ Lattice: Radial Topology, Zeta Holonomy, and the Algebraic Symmetrieleiter in the EABC Framework (v4)}}
\author{\textit{Dr. Thomas Hoffbauer} \\ \small EABC - (E)pstein-Zeta-based (A)ithmetic (B)oundary-condition and (C)oarse-graining) framework -  Bamberg }
\date{August 16, 2026}

\begin{document}
	
	\maketitle
	
	\begin{abstract}
		We present a unified mathematical-physical formulation of the EABC (Epstein-Zeta-based Arithmetic Boundary-condition and Coarse-graining) framework. Operating as a structured hierarchy of algebraic and geometric spaces, the framework links the discrete distribution of prime numbers to the exceptional Lie group $E_8$. Starting with Arithmetic Supersymmetry (ArithSUSY) on square-free integers, we ascend the algebraic ``Symmetrieleiter'' through Lipschitz-quaternion norm-shells, non-associative octonionic structures ($G_2$), exceptional Jordan algebras ($E_6$), Freudenthal triple systems ($E_7$), and ultimately the universally optimal $E_8$ lattice. By introducing a systematic coarse-graining projection of number-theoretic microstates into support-class Markov chains, we demonstrate the emergence of detailed balance, local temporal persistence, and study a heuristic model of chiral phase dynamics without violating the underlying time-reversal symmetry of the arithmetic vacuum. Rigorous epistemic demarcation matrices (Ebene [C] Claim-Walls) are enforced at each stage to preserve the formal boundary between verified mathematical theorems and heuristic physical analogies.
	\end{abstract}
	
	\tableofcontents
	\newpage
	
	
	\section{Arithmetic Supersymmetry (ArithSUSY) on the Square-Free Lattice}
	\label{sec:arith_susy}
	
	We formalize the algebraic structure of Arithmetic Supersymmetry (ArithSUSY) within the square-free sector of the natural numbers $\mathbb{N}_{\text{sf}}$. This formulation establishes a rigorous isomorphism between multiplicative number theory and a 1D discrete supersymmetric quantum field theory, wherein the prime factors of an integer act as independent fermionic mode excitations.
	
	\subsection{The Fermionic Parity Operator and State Space}
	Let $\mathbb{N}_{\text{sf}}$ denote the set of square-free positive integers. For any $n \in \mathbb{N}_{\text{sf}}$, the number of distinct prime factors is given by the additive function $\omega(n)$. We define the fermionic Fock space $\mathcal{H}_{\text{sf}} = \ell^2(\mathbb{N}_{\text{sf}})$, which decomposes into orthogonal bosonic and fermionic subspaces:
	\begin{equation}
		\mathcal{H}_{\text{sf}} = \mathcal{H}_{\text{B}} \oplus \mathcal{H}_{\text{F}},
	\end{equation}
	where the subspaces are defined according to the parity of the prime factor count:
	\begin{align}
		\mathcal{H}_{\text{B}} &= \operatorname{span} \left\{ |n\rangle : n \in \mathbb{N}_{\text{sf}}, \, \omega(n) \equiv 0 \pmod 2 \right\}, \\
		\mathcal{H}_{\text{F}} &= \operatorname{span} \left\{ |n\rangle : n \in \mathbb{N}_{\text{sf}}, \, \omega(n) \equiv 1 \pmod 2 \right\}.
	\end{align}
	The canonical fermion parity operator $(-1)^F$ on $\mathcal{H}_{\text{sf}}$ is represented exactly by the classical Möbius function $\mu(n)$ for square-free inputs:
	\begin{equation}
		\langle n | (-1)^F | n \rangle = \mu(n) = (-1)^{\omega(n)}.
	\end{equation}
	This identification establishes that the Möbius parity acts as the exact arithmetic equivalent of the physical $(-1)^F$ operator.
	
	\subsection{Asymptotic 50/50 Balance and the Supertrace-Structure}
	A fundamental requirement of unbroken supersymmetry is the exact degeneracy between bosonic and fermionic states at non-zero energy levels. In our multiplicative number theory framework, this balance is verified via the asymptotic distribution of odd and even square-free numbers. Let $Q(N)$ be the number of square-free integers up to $N$, which is asymptotically given by:
	\begin{equation}
		Q(N) = \sum_{n \le N} |\mu(n)| = \frac{6}{\pi^2} N + \mathcal{O}(\sqrt{N}).
	\end{equation}
	Let $Q_{\text{B}}(N)$ and $Q_{\text{F}}(N)$ denote the count of bosonic ($\mu(n) = 1$) and fermionic ($\mu(n) = -1$) states, respectively. The cumulative Mertens function $M(N)$ measures the exact state imbalance:
	\begin{equation}
		M(N) = \sum_{n \le N} \mu(n) = Q_{\text{B}}(N) - Q_{\text{F}}(N).
	\end{equation}
	Numerical audits up to $N = 10^6$ establish a state ratio of:
	\begin{equation}
		\frac{Q_{\text{B}}(10^6)}{Q(10^6)} \approx 50.02\%, \quad \frac{Q_{\text{F}}(10^6)}{Q(10^6)} \approx 49.98\%,
	\end{equation}
	demonstrating the extremely rapid asymptotic convergence toward a perfect 50/50 supersymmetric balance in the arithmetic vacuum.
	
	This algebraic balance is beautifully synthesized by the canonical \textbf{Arithmetic Supertrace Identity} on the square-free Fock sector. For the operator $H|n\rangle = (\log n)|n\rangle$, the partition function is exactly the inverse of the Riemann Zeta function:
	\begin{equation}
		\label{eq:supertrace_identity}
		\operatorname{Tr}\left[(-1)^F e^{-\beta H}\right] = \sum_{\substack{n \ge 1 \\ n \text{ squarefree}}} \frac{\mu(n)}{n^\beta} = \frac{1}{\zeta(\beta)} \quad (\Re \beta > 1).
	\end{equation}
	This supertrace provides a rigorous formal mapping between the fermionic path integral of a 1D supersymmetric field theory and the poles and zeros of the Riemann Zeta function, establishing an algebraic foundation for the vacuum state.
	
	\subsection{The Mertens Bound, Riemann Hypothesis, and the historical Odlyzko--te Riele Stress-Gate}
	The fluctuations of the Mertens function $M(N)$ characterize the stability of the ArithSUSY vacuum. Under the Riemann Hypothesis (RH), the growth of the state imbalance is strictly constrained by the spectral properties of the non-trivial Zeta zeros:
	\begin{equation}
		M(N) = \mathcal{O}\left(N^{\frac{1}{2} + \varepsilon}\right), \quad \forall \varepsilon > 0.
	\end{equation}
	This is equivalent to the assertion that the non-trivial zeros lie exactly on the critical line $\Re(s) = 1/2$. In our numerical audit up to $N = 10^6$, the absolute state deviation is bounded by:
	\begin{equation}
		\max_{N \le 10^6} |M(N)| = 368,
	\end{equation}
	which is tightly constrained well below the asymptotic SUSY-breaking threshold of $\sqrt{N} = 1000$.
	
	Historically, Franz Mertens conjectured in 1897 that $|M(N)| < \sqrt{N}$ for all $N > 1$. This strong version of the bound was spectacularly disproved in 1985 by Andrew Odlyzko and Herman te Riele, who showed that:
	\begin{equation}
		\limsup_{N \to \infty} \frac{M(N)}{\sqrt{N}} > 1.06 \quad \text{and} \quad \text{liminf}_{N \to \infty} \frac{M(N)}{\sqrt{N}} < -1.009.
	\end{equation}
	Within our scientific governance, this disproof is not treated as a failure or a structural weakness of the model. Rather, it serves as a \textbf{perfect Stress-Gate for the ArithSUSY vacuum}. The Odlyzko--te Riele result proves that while the spectral properties of the non-trivial zeros constrain the fluctuations on average, they allow for extremely rare, non-local constructive interferences where the vacuum \"jumps\" beyond the naive $\pm \sqrt{N}$ boundary.
	
	\subsection{Mertens Fluctuation Audit against Order-Matched Null Models: Empirical Findings}
	To verify whether the fluctuations observed in the Mertens function exhibit any genuine arithmetic features of ArithSUSY or simply statistical properties of local correlations, we subject the Möbius sequence $\mu(n)$ (restricted to square-free numbers) to a rigorous comparison against three distinct classes of null models:
	\begin{enumerate}
		\item $\mathbf{N_0}$ (Shuffle Model): Completely destroys all sequential dependencies through random permutation of the Möbius sequence.
		\item $\mathbf{N_1}$ (Markov-1 Model): Preserves the first-order transition probabilities $P(\mu(n+1) | \mu(n))$ but discards all higher-order memory.
		\item $\mathbf{N_2}$ (Markov-2 Model): Preserves the second-order transition probabilities $P(\mu(n+2) | \mu(n+1), \mu(n))$.
	\end{enumerate}
	
	To evaluate these models, we define four independent statistical observables:
	\begin{enumerate}
		\item \textbf{Scaled Maximum Statistic ($R_N$):} Measures the peak normalized deviation up to scale $N$:
		\begin{equation}
			R_N = \max_{n \le N} \frac{|M(n)|}{\sqrt{n}}.
		\end{equation}
		\item \textbf{Variance of the Cumulative Mean ($V(N)$):}
		\begin{equation}
			V(N) = \frac{1}{N} \sum_{n \le N} M(n)^2.
		\end{equation}
		\item \textbf{Block Increments ($\Delta_L M(n)$):} For a block length $L$:
		\begin{equation}
			\Delta_L M(n) = M(n+L) - M(n) = \sum_{k=1}^L \mu(n+k).
		\end{equation}
		\item \textbf{Autocorrelation ($C(k)$):} Measures the correlation at lag $k$:
		\begin{equation}
			C(k) = \frac{1}{N} \sum_{n \le N} \mu(n)\mu(n+k).
		\end{equation}
	\end{enumerate}
	
	Our quantitative audit of the Möbius sequence up to $N = 10^6$ reveals a fundamental two-fold division:
	\begin{itemize}
		\item \textbf{Local Correlations ($C(k)$ and Block Increments $\Delta_L M$):} These are completely absorbed by the Markov-1 and Markov-2 null models ($N_1, N_2$) as well as the exact shuffle, yielding $p$-values between $0.21$ and $0.96$. On a local scale, there is no statistically significant evidence of unexplained dynamical rigidity.
		\item \textbf{Global Cumulative Amplitudes ($R_N$ and $V(N)$):} The large-$N$ maximal statistic ($R_{10^3} \approx 0.47220, R_{10^4} \approx 0.46298$) and the cumulative variance ($V(10^6) = 14851.09$) lie significantly at the lower boundary compared to the count-preserving shuffle ($p_{\text{MC,two}} = 0.0165$). The sequence thus exhibits unusually small global fluctuations up to $10^6$.
	\end{itemize}
	Since the standard $N_1$ and $N_2$ surrogates do not exactly preserve the global count boundary of the frequencies $\#\{-1\}, \#\{0\}, \#\{+1\}$, it remains an important open gate (\emph{Open Gate}) whether this global signal can be completely resolved by an exact count-conditioned Markov-2 specification, or if it persists as a genuine long-range structural feature.
	
	\subsection{Formal Demarcation (AS.5 KERNEL: Modul-2 Audit-Fassung)}
	To maintain the absolute formal integrity of the EABC repository, we enforce a strict demarcation boundary regarding the physical interpretation of ArithSUSY and Mertens fluctuations, freezing the final audited status in Version v4.7:
	
	\begin{equation}
		\begin{array}{l|c|l}
			\hline
			\textbf{Observable / Hypothese} & \textbf{Status} & \textbf{Epistemische Kategorie} \\
			\hline
			\text{Lokale } C(k) \text{-Struktur bis Lag 20} & \mathrm{ABSORBED} & \text{Ebene A: Markov-reproduzierbar} \\
			\Delta_L M \text{ für } L = 10, 100, 1000 & \mathrm{ABSORBED} & \text{Ebene A: Markov-reproduzierbar} \\
			R_N \text{ gegen exakten Shuffle} & \mathrm{SIGNAL} & \text{Ebene A: Signifikante Abweichung} \\
			V(N) \text{ gegen exakten Shuffle} & \mathrm{SIGNAL} & \text{Ebene A: Signifikante Abweichung} \\
			R_N, V(N) \text{ jenseits count-gematchtem } N_2 & \mathrm{OPEN} & \text{Ebene B: Offenes Forschungs-Gate} \\
			\text{GUE als physikalische Ursache} & \mathrm{NOT\ PROVEN} & \text{Ebene C: Unbewiesene Kausal-Zuweisung} \\
			\text{Physische ArithSUSY-Stabilisierung} & \mathrm{NOT\ CLAIMED} & \text{Ebene C: Heuristische Analogie} \\
			\hline
		\end{array}
	\end{equation}
	While the supertrace mapping and the Fock space structures are exact mathematical theorems (Ebene A), the interpretation of GUE-stiffness as a physical "stabilizing pressure" or ArithSUSY as a real spacetime physical mechanism is strictly a heuristic model and carries no physical spacetime causal claim.
	
	\subsection{Thermal Parity Crossover and Analytical Partition Functions}
	We study the thermodynamic behavior of the ArithSUSY state space under a temperature-weighted projection, defining the partition functions for the bosonic and fermionic sectors. Rather than relying solely on numerical approximations, we show that the partition functions are analytically solvable in closed form.
	
	For any inverse temperature $\beta > 1$ (the scaling dimension of the Fock space), the total partition function of the square-free sector $\mathcal{H}_{\text{sf}}$ is exactly given by the Dirichlet series:
	\begin{equation}
		Z_{\text{tot}}(\beta) = \sum_{n \in \mathbb{N}_{\text{sf}}} n^{-\beta} = \sum_{n \ge 1} \frac{|\mu(n)|}{n^\beta} = \frac{\zeta(\beta)}{\zeta(2\beta)},
	\end{equation}
	while the superpartition function (or supersymmetric supertrace) is represented by:
	\begin{equation}
		Z_{\text{super}}(\beta) = \operatorname{Tr}\left[(-1)^F e^{-\beta H}\right] = \sum_{n \ge 1} \frac{\mu(n)}{n^\beta} = \frac{1}{\zeta(\beta)}.
	\end{equation}
	
	By combining these two relations, the bosonic partition function $Z_{\text{B}}(\beta) = \frac{1}{2}(Z_{\text{tot}} + Z_{\text{super}})$ and the fermionic partition function $Z_{\text{F}}(\beta) = \frac{1}{2}(Z_{\text{tot}} - Z_{\text{super}})$ can be expressed in exact analytical closed form:
	\begin{equation}
		\boxed{
			Z_{\text{B}}(\beta) = \frac{1}{2} \left[ \frac{\zeta(\beta)}{\zeta(2\beta)} + \frac{1}{\zeta(\beta)} \right]
		}
	\end{equation}
	and
	\begin{equation}
		\boxed{
			Z_{\text{F}}(\beta) = \frac{1}{2} \left[ \frac{\zeta(\beta)}{\zeta(2\beta)} - \frac{1}{\zeta(\beta)} \right].
		}
	\end{equation}
	
	From these closed forms, several important analytical limits follow directly:
	\begin{enumerate}
		\item \textbf{The High-Temperature Limit ($\beta \to 1^+$):} As the temperature approaches the boundary of the pole of the Riemann Zeta function, $Z_{\text{tot}}(\beta) \to \infty$ while $Z_{\text{super}}(\beta) \to 0$. Consequently, the temperature-weighted parity ratios satisfy:
		\begin{equation}
			\lim_{\beta \to 1^+} \frac{Z_{\text{B}}(\beta)}{Z_{\text{tot}}(\beta)} = \frac{1}{2}, \quad \text{and} \quad \lim_{\beta \to 1^+} \frac{Z_{\text{F}}(\beta)}{Z_{\text{tot}}(\beta)} = \frac{1}{2}.
		\end{equation}
		This asymptotic 50/50 balance is an exact analytical theorem, rather than a phenomenological or dynamical "restoration" of supersymmetry.
		\item \textbf{The Low-Temperature Limit ($\beta \to \infty$):} As the scale parameter becomes arbitrarily large, the partition functions collapse to:
		\begin{equation}
			\lim_{\beta \to \infty} Z_{\text{B}}(\beta) = 1, \quad \text{and} \quad \lim_{\beta \to \infty} Z_{\text{F}}(\beta) = 0.
		\end{equation}
		This corresponds to the fact that $n=1$ is the unique ground state of the system and carries bosonic parity ($\mu(1)=1$), while the lowest-energy fermionic state is $n=2$ ($\mu(2)=-1$). Indeed, the fermionic states sum over primes and their odd square-free products, while bosonic states sum over the identity and even products:
		\begin{align}
			Z_{\text{F}}(\beta) &= 2^{-\beta} + 3^{-\beta} + 5^{-\beta} + \dots, \\
			Z_{\text{B}}(\beta) &= 1 + 6^{-\beta} + 10^{-\beta} + 14^{-\beta} + \dots.
		\end{align}
	\end{enumerate}
	
	To maintain rigorous physical discipline, the continuous $\beta$-dependent separation of parity states must be classified as a smooth \textbf{thermal parity crossover} or \textbf{temperature-weighted parity separation}, rather than a true thermodynamic phase transition, as the partition functions $Z_{\text{B}}(\beta)$ and $Z_{\text{F}}(\beta)$ remain strictly analytic for all $\beta > 1$.
	
	Numerical verification up to $N=10^6$ confirms this analytical structure. For instance, at $\beta = 3.04$, the numerical audit yields $Z_{\text{super}}(3.04) = Z_{\text{B}}(3.04) - Z_{\text{F}}(3.04) \approx 0.837216$, matching the exact theoretical value of $1/\zeta(3.04) \approx 0.83721639$ to within machine precision without requiring any non-local parameter adjustments. We establish this analytical crossover within the formal Claim-Wall:
	
	\begin{equation}
		\begin{array}{l|c}
			\hline
			\textbf{Aussage / Hypothese} & \textbf{Sicherheits-Status} \\
			\hline
			Z_{\rm super}=1/\zeta(\beta) & \mathrm{PROVEN} \\
			Z_{\rm tot}=\zeta(\beta)/\zeta(2\beta) & \mathrm{PROVEN} \\
			\text{Geschlossene Formeln f"ur } Z_B, Z_F & \mathrm{PROVEN} \\
			B/F \to 50/50 \text{ f"ur } \beta \to 1^+ & \mathrm{PROVEN\ LIMIT} \\
			Z_B \to 1, \; Z_F \to 0 \text{ f"ur } \beta \to \infty & \mathrm{PROVEN\ LIMIT} \\
			\text{Thermischer Crossover} & \mathrm{MATHEMATICAL\ DESCRIPTION} \\
			\text{Echter Phasen"ubergang} & \mathrm{NOT\ ESTABLISHED} \\
			\text{Physikalische SUSY} & \mathrm{NOT\ CLAIMED} \\
			\hline
		\end{array}
	\end{equation}
	
	\subsection{Finite Group Fourier Decomposition modulo 3: Channel Separation}
	\label{subsec:mod3_fourier}
	We formalize the orthogonal decomposition of state-space amplitudes over the reduced residue group modulo $3$, $(\mathbb{Z}/3\mathbb{Z})^\times = \{1, 2\} \cong C_2$, demonstrating that the factor of $3^2 = 9$ appearing in the bare coupling definition represents a strictly defined algebraic dimension of the dual tensor product of these channels, rather than a physical spatial dimensional claim.
	
	The group $(\mathbb{Z}/3\mathbb{Z})^\times$ possesses a single non-trivial quadratic Dirichlet character (the Legendre symbol $\chi_1(n) = \left(\frac{n}{3}\right)$) and the trivial principal character $\chi_0(n)$. Any complex-valued amplitude function $f = (f(1), f(2))^T \in \ell^2((\mathbb{Z}/3\mathbb{Z})^\times)$ can be decomposed orthogonally into symmetric and antisymmetric Fourier sectors under the action of the cyclic group $C_2$.
	
	We define the unitary discrete Fourier transform (or Hadamard-Walsh transformation) $U$ acting on the amplitude vector:
	\begin{equation}
		U = \frac{1}{\sqrt{2}} \begin{pmatrix} 1 & 1 \\ 1 & -1 \end{pmatrix}.
	\end{equation}
	Applying $U$ to the amplitude vector $f$ projects the degrees of freedom onto two decoupled channels:
	\begin{align}
		\widehat{f}_0 &= \frac{f(1) + f(2)}{\sqrt{2}} \quad \text{(Symmetric Channel)}, \\
		\widehat{f}_1 &= \frac{f(1) - f(2)}{\sqrt{2}} \quad \text{(Antisymmetric Channel)}.
	\end{align}
	The orthogonality of these channels is strictly guaranteed by the Schur relations:
	\begin{equation}
		\langle \widehat{f}_0, \widehat{f}_1 \rangle = 0 \quad \forall f \in \ell^2((\mathbb{Z}/3\mathbb{Z})^\times).
	\end{equation}
	
	When evaluating the dual interaction of these channels in the bilinear coupling, the tensor product space $\mathcal{H}_{\text{coupling}} = \ell^2((\mathbb{Z}/3\mathbb{Z})^\times) \otimes \ell^2((\mathbb{Z}/3\mathbb{Z})^\times)$ spans exactly a 4-dimensional space of discrete coupling coefficients, which represents the discrete algebraic core of the regularized 3-dimensional Euclidean metric tensor space $\mathbb{R}^3 \otimes \mathbb{R}^3$ of dimension $d^2 = 3^2 = 9$. We freeze the formal demarcation of this structure:
	
	\begin{equation}
		\begin{array}{l|c|l}
			\hline
			\textbf{Aussage / Relation} & \textbf{Sicherheits-Status} & \textbf{Epistemische Kategorie} \\
			\hline
			\text{Charakter-Orthogonalit\"{a}t modulo } 3 & \mathrm{PROVEN} & \text{Ebene A: Zahlentheoretischer Satz} \\
			\text{Orthogonale } C_2\text{-Fourier-Zerlegung via } U & \mathrm{PROVEN} & \text{Ebene A: Unit\"{a}re Lie-Theorie} \\
			\dim(\mathbb{R}^3 \otimes \mathbb{R}^3) = 9 & \mathrm{PROVEN} & \text{Ebene A: Lineare Algebra} \\
			\{+1, -1, 0\} \leftrightarrow \text{Raum-Orientierungen} & \mathrm{HEURISTIC} & \text{Ebene C: Strukturanalogie} \\
			3^2 \text{ in } \alpha_{\text{bare}} \text{ entspringt mod-3-Fourier} & \mathrm{MODEL\ ASSUMPTION} & \text{Ebene C: Regularisierungs-Eichung} \\
			\text{Ableitung der realen QED-Kopplungsst\"{a}rke} & \mathrm{NOT\ ESTABLISHED} & \text{Ebene C: Keine exp. Signatur} \\
			\hline
		\end{array}
	\end{equation}
	This rigorous channel separation protects the framework from any naive spatial overinterpretations, defining the modulo-3 factor as an exact algebraic dimension of the dual Fourier state space.
	
	\newpage
	
	\section{Epstein-Zeta Holonomy and the Topological Phase Collapse}
	\label{sec:epstein_zeta_holonomy}
	
	We study the local topological phase structure of the unperturbed $E_8$ lattice vacuum through the holomorphic flow of the radial Epstein-Zeta function $Z^{\mathrm{len}}_{E_8}(s)$ in the local regions $0 < R < 10$ and $10 < R < 12$. At the boundary of the near spectral region, the dynamics is governed by the meromorphy of the spectral divisor, which physically acts as a coupled pair of topological defects. We prove that the resulting $U(1)$ winding behavior exhibits an abrupt phase transition driven by exact charge compensation, including the trivial zero compensation at $s = 2$.
	
	\subsection{The Meromorphic Two-Defect Model}
	In the spectral representation, the unperturbed $E_8$ lattice state-sum factors into a product of shifted Riemann Zeta functions:
	\begin{equation}
		Z^{\mathrm{len}}_{E_8}(s) = 240 \cdot 2^{-s/2} \cdot \zeta\left(\frac{s}{2}\right) \cdot \zeta\left(\frac{s}{2} - 3\right).
	\end{equation}
	From the logarithmic derivative, we define the closed meromorphic 1-form:
	\begin{equation}
		\omega_Z(s) = d\log Z^{\mathrm{len}}_{E_8}(s) = \frac{Z'_{E_8}(s)}{Z_{E_8}(s)} \, ds.
	\end{equation}
	In the local spectral vicinity, this differential form is modeled by two entropic defect centers:
	\begin{enumerate}
		\item \textbf{The Pole (Anti-Vortex, charge $q_{\mathrm{pole}} = -1$):} Positioned at $s_{\mathrm{pole}} = 8$, representing the embedding dimension of the $E_8$ lattice.
		\item \textbf{The Zero (Vortex, charge $q_{\mathrm{zero}} = +1$):} Positioned at $s_{\mathrm{zero}} = -2$, originating from the first trivial zero of the shifted zeta factor.
	\end{enumerate}
	The local approximation of the 1-form is given by:
	\begin{equation}
		\omega_Z(s) \approx -\frac{1}{s-8} + \frac{1}{s+2}.
	\end{equation}
	
	\subsection{Singular Contour Stratification and Winding Numbers}
	We evaluate the winding number $\nu(R)$ along a circular contour $\gamma(t) = 8 + R e^{i t}$ ($t \in [0, 2\pi]$) centered at the pole $s = 8$. By the argument principle of complex analysis, the winding number is a piecewise constant step function with a singular jump at the critical boundary $R_0 = 10$:
	\begin{equation}
		\nu(R) = \frac{1}{2\pi i} \oint_{\gamma} \omega_Z(s) \, ds = 
		\begin{cases} 
			-1, & 0 < R < 6, \\
			0, & 6 < R < 10, \\
			1, & 10 < R < 12.
		\end{cases}
	\end{equation}
	
	The transition radii $R_0 = 6$ and $R_1 = 10$ are rigidly determined by the Euclidean distances between the pole and the zeros:
	\begin{equation}
		R_0 = |8 - 2| = 6, \quad R_1 = |8 - (-2)| = 10.
	\end{equation}
	
	For $0 < R < 6$, only the pole is enclosed within the contour ($N_Z = 0, N_P = 1$), yielding $\nu = -1$. For $6 < R < 10$, the pole and the first zero at $s=2$ are enclosed ($N_Z = 1, N_P = 1$), leading to a complete charge neutralization ($\nu = 0$). For $10 < R < 12$, the pole and both zeros at $s=2, -2$ are enclosed, giving $\nu = +1$.
	
	\subsection{Topological Phase Collapse and Holonomy Behavior}
	We define the topological phase collapse as the discrete loss of the net winding number upon enclosing a second defect of opposite topological charge. The transition is structurally classified as:
	\begin{equation}
		\text{Pole Sector } (\nu = -1) \quad \xrightarrow{R \nearrow 6} \quad \text{First Zero Compensation } (\nu = 0) \quad \xrightarrow{R \nearrow 10} \quad \text{Second Zero Contribution } (\nu = 1).
	\end{equation}
	In terms of effective topological charges, these transitions represent the sequential enclosing of vortices:
	\begin{equation}
		-1 \quad \longrightarrow \quad (-1) + (+1) = 0 \quad \longrightarrow \quad 0 + (+1) = 1.
	\end{equation}
	
	The actual topological observable is the unexponentiated winding number itself, obtained by integrating the 1-form along the loop $\gamma$:
	\begin{equation}
		\nu(R) = \frac{1}{2\pi i} \oint_{\gamma} \omega_Z(s) \, ds.
	\end{equation}
	This winding number directly detects the topological phase of the divisor. Unlike an exponentiated holonomy of the form $e^{2\pi i \nu}$ (which trivially reduces to $1$ for all integer winding sectors and thus washes out the transition), the winding number acts as a sharp topological indicator across the local transitions. The strong numerical amplitude fluctuations observed near the critical boundaries are a consequence of the singular contour contact with the zeros at $s = 2$ and $s = -2$, representing expected singular behavior at the boundaries.
	
	\subsection{Epistemic Demarcation of the Phase Boundary}
	To maintain the formal security and mathematical precision of the EABC repository, we enforce a strict demarcation boundary regarding the physical interpretation of the spectral phase transition:
	
	\begin{equation}
		\begin{array}{l|c}
			\textbf{Structure / Parameter} & \textbf{Epistemic Status (Governance Category)} \\
			\hline
			s_{\rm pole}=8,\ s_{\rm zero1}=2,\ s_{\rm zero2}=-2 & \mathrm{DEFINED\ (Ebene\ A)} \\
			R_0=6,\ R_1=10 & \mathrm{EXACT\ (Ebene\ A)} \\
			\nu(0 < R < 6)=-1 & \mathrm{EXACT\ (Ebene\ A)} \\
			\nu(6 < R < 10)=0 & \mathrm{EXACT\ (Ebene\ A)} \\
			\nu(10 < R < 12)=1 & \mathrm{EXACT\ (Ebene\ A)} \\
			\text{Winding number transitions} & \mathrm{EXPECTED\ SINGULAR\ BEHAVIOR\ (Ebene\ A)} \\
			\text{Physical phase transition in spacetime} & \mathrm{NOT\ CLAIMED\ (Ebene\ C)} \\
		\end{array}
	\end{equation}
	
	This rigorous classification shields the EABC-AS.5 kernel from unwarranted physical extrapolation, while preserving the exact, mathematically proven topological properties of the Epstein-Zeta phase boundary.
	
	
	\section{Arithmetic Stratification of Prime Quaternion Shells}
	\label{sec:prime_shell_stratification}
	
	We study the arithmetic and geometric structures of prime norm shells in the Lipschitz quaternion lattice. The primary goal of this section is to prove that the minimal support of elements in these shells is entirely governed by prime congruences modulo $24$. This algebraic stratification induces a rigorous stabilizer and orbit hierarchy under the action of the hyperoctahedral group, establishing a coordinate-free bridge from modular arithmetic to discrete geometry.
	
	\subsection{Prime norm shells in the Lipschitz lattice}
	Let $\mathbb{H}_{\mathbb{Z}}$ denote the ring of Lipschitz integers (also known as the Lipschitz quaternion lattice), defined as:
	\begin{equation}
		\mathbb{H}_{\mathbb{Z}} = \left\{ q = x_1 + x_2 i + x_3 j + x_4 k : x_j \in \mathbb{Z} \right\}.
	\end{equation}
	The norm of a Lipschitz quaternion $q \in \mathbb{H}_{\mathbb{Z}}$ is given by the positive definite quadratic form $N(q) = q\bar{q} = x_1^2 + x_2^2 + x_3^2 + x_4^2$. For a prime number $p > 3$, the prime norm shell $S_p$ is defined as the set of all Lipschitz integers of norm $p$:
	\begin{equation}
		S_p = \left\{ x \in \mathbb{Z}^4 : x_1^2 + x_2^2 + x_3^2 + x_4^2 = p \right\}.
	\end{equation}
	By Jacobi's four-square theorem, the cardinality of $S_p$ for any odd prime $p$ is exactly $|S_p| = 8\sigma_1(p) = 8(p+1)$, indicating that these shells are always non-empty and grow linearly with $p$.
	
	\subsection{Minimal support and congruence classes}
	To analyze the geometric distribution of points within $S_p$, we introduce the concept of the support of a vector. For any $x = (x_1, x_2, x_3, x_4) \in S_p$, the support $\operatorname{supp}(x)$ is the set of indices with non-zero coordinates:
	\begin{equation}
		\operatorname{supp}(x) = \left\{ j \in \{1, 2, 3, 4\} : x_j \neq 0 \right\}.
	\end{equation}
	The support size is denoted by $s(x) = \#\operatorname{supp}(x)$. We define the minimal support of the prime shell $S_p$ as the minimum support size over all points in the shell:
	\begin{equation}
		s_{\min}(p) = \min_{x \in S_p} s(x) = \min_{x \in S_p} \#\left\{ j : x_j \neq 0 \right\}.
	\end{equation}
	Geometrically, $s_{\min}(p)$ determines the lowest-dimensional coordinate subspace of $\mathbb{R}^4$ that intersects the prime shell $S_p$. If $s_{\min}(p) = 2$, the shell contains points lying on 2-dimensional planes spanned by coordinate axes (e.g., of the form $(a, b, 0, 0)$). If $s_{\min}(p) = 3$, the lowest-dimensional intersection occurs on 3-dimensional coordinate hyperplanes (e.g., $(a, b, c, 0)$). If $s_{\min}(p) = 4$, points in $S_p$ are strictly excluded from all coordinate hyperplanes, meaning every coordinate of every vector in $S_p$ is non-zero.
	
	\subsection{The mod-$24$ support theorem}
	We now prove the central algebraic result, which shows that the minimal support $s_{\min}(p)$ of a prime shell is a rigid arithmetic invariant determined entirely by the residue of $p$ modulo $24$.
	
	\begin{theorem}[Modulo-24 Support Stratification]
		\label{thm:mod24_support}
		Let $p > 3$ be a prime number. The minimal support $s_{\min}(p)$ in the Lipschitz shell $S_p$ is given exactly by:
		\begin{equation}
			s_{\min}(p) = 
			\begin{cases} 
				2, & p \bmod 24 \in \{1, 5, 13, 17\}, \\ 
				3, & p \bmod 24 \in \{11, 19\}, \\ 
				4, & p \bmod 24 \in \{7, 23\}. 
			\end{cases}
		\end{equation}
	\end{theorem}
	
	\begin{proof}
		The proof proceeds by analyzing the solubility of the sum-of-squares relation in fewer coordinates, utilizing classical results of Fermat, Legendre--Gauss, and Lagrange.
		
		\noindent \textbf{Case 1: $s_{\min}(p) = 2$.}
		A support size of $2$ is achievable if and only if there exist integers $x_1, x_2 \in \mathbb{Z} \setminus \{0\}$ such that $x_1^2 + x_2^2 = p$. By Fermat's two-square theorem, a prime $p$ can be written as the sum of two squares if and only if $p \equiv 1 \pmod 4$. Since $p > 3$, we analyze this condition modulo $24$:
		\begin{equation}
			p \equiv 1 \pmod 4 \iff p \bmod 24 \in \{1, 5, 13, 17\}.
		\end{equation}
		Thus, for any prime in these four residue classes, there exists a representation of the form $x_1^2 + x_2^2 = p$, and since $p$ is prime, neither coordinate can be zero. Thus, $s_{\min}(p) = 2$ in these cases.
		
		\noindent \textbf{Case 2: $s_{\min}(p) = 3$.}
		A support size of $3$ is achievable if and only if $p$ can be written as a sum of three squares, but not as a sum of two squares. This requires $p$ to be representable as $x_1^2 + x_2^2 + x_3^2 = p$ with non-zero coordinates. By Legendre's three-square theorem, an integer $n$ can be written as the sum of three squares if and only if $n$ is not of the form $4^a(8b + 7)$. For an odd prime $p > 3$, this condition simplifies to $p \not\equiv 7 \pmod 8$. 
		
		Analyzing this modulo $24$, the primes that can be represented as the sum of three squares must satisfy:
		\begin{equation}
			p \bmod 24 \in \{1, 5, 11, 13, 17, 19\}.
		\end{equation}
		To find the primes where the \emph{minimal} support is exactly $3$, we must exclude the primes that already admit a two-square representation (Case 1). Subtracting $\{1, 5, 13, 17\}$ from the set of three-square representable primes leaves exactly:
		\begin{equation}
			p \bmod 24 \in \{11, 19\}.
		\end{equation}
		For these classes, $p$ can be written as $x_1^2 + x_2^2 + x_3^2 = p$, but $p \not\equiv 1 \pmod 4$ prevents any two-coordinate representation. Thus, $s_{\min}(p) = 3$ for these residue classes.
		
		\noindent \textbf{Case 3: $s_{\min}(p) = 4$.}
		By Lagrange's four-square theorem, every integer can be represented as the sum of four squares. Therefore, any prime $p$ admits a representation $x_1^2 + x_2^2 + x_3^2 + x_4^2 = p$. The minimal support is exactly $4$ if and only if $p$ cannot be represented as a sum of three or fewer squares. 
		
		According to Legendre's theorem, this occurs if and only if $p \equiv 7 \pmod 8$. Expressing this modulo $24$ yields:
		\begin{equation}
			p \equiv 7 \pmod 8 \iff p \bmod 24 \in \{7, 23\}.
		\end{equation}
		For these primes, any representation in $S_p$ must have all four coordinates non-zero, ensuring that $s_{\min}(p) = 4$. This completes the proof of Theorem~\ref{thm:mod24_support}.
	\end{proof}
	
	\subsection{Hyperoctahedral symmetry and stabilizers}
	The natural symmetry group acting on the Lipschitz shell $S_p$ is the hyperoctahedral group $B_4$, which has order 384, in contrast to the Weyl group of the $D_4$ root system which has order 192. It is defined as the semidirect product:
	\begin{equation}
		B_4 = (\mathbb{Z}/2\mathbb{Z})^4 \rtimes S_4,
	\end{equation}
	where $(\mathbb{Z}/2\mathbb{Z})^4$ acts by independent coordinate sign changes and the symmetric group $S_4$ acts by permutation of the four coordinate axes. The order of the group is $|B_4| = 2^4 \cdot 4! = 384$.
	
	The stratification of $S_p$ by support size induces a corresponding stratification of the stabilizers under the $B_4$ action. Let $x \in S_p$ be a vector, and let $\operatorname{Stab}_{B_4}(x) = \{ g \in B_4 : g \cdot x = x \}$ be its stabilizer subgroup. The structure of this stabilizer is determined by the number of zero coordinates and the multiplicity of identical absolute values among the non-zero coordinates.
	
	\subsection{Generic orbit cardinalities}
	We establish the exact orbit size hierarchy under the action of $B_4$. For any point $x \in S_p$, the orbit is denoted by $O(x) = \{ g \cdot x : g \in B_4 \}$. By the orbit-stabilizer theorem, the cardinality of the orbit is:
	\begin{equation}
		|O(x)| = \frac{|B_4|}{|\operatorname{Stab}_{B_4}(x)|} = \frac{384}{|\operatorname{Stab}_{B_4}(x)|}.
	\end{equation}
	
	To identify the generic or maximal orbit size for each support stratum, we assume that the non-zero coordinates have pairwise distinct positive absolute values. Under this generic assumption, the stabilizer structures and orbit sizes are classified as follows:
	
	\begin{theorem}[Generic Orbit Size Hierarchy]
		\label{thm:generic_orbits}
		Let $x \in S_p$ be a vector belonging to a minimal-support stratum of size $s(x) = s_{\min}(p)$. Assuming that the active non-zero coordinates have distinct positive absolute values, the stabilizer subgroups and orbit sizes are given by:
		\begin{enumerate}
			\item \textbf{Support $2$ (Generic form $x = (a, b, 0, 0)$ with $0 < a < b$):}
			The inactive coordinates ($x_3, x_4 = 0$) can be signed and permuted freely among themselves, yielding a stabilizer isomorphic to $B_2 \cong (\mathbb{Z}/2\mathbb{Z})^2 \rtimes S_2$ of order $8$.
			\begin{equation}
				|\operatorname{Stab}_{B_4}(x)| = 8 \implies |O(x)| = \frac{384}{8} = 48.
			\end{equation}
			
			\item \textbf{Support $3$ (Generic form $x = (a, b, c, 0)$ with $0 < a < b < c$):}
			The single inactive coordinate ($x_4 = 0$) can have its sign flipped freely, leading to a stabilizer isomorphic to $B_1 \cong \mathbb{Z}/2\mathbb{Z}$ of order $2$.
			\begin{equation}
				|\operatorname{Stab}_{B_4}(x)| = 2 \implies |O(x)| = \frac{384}{2} = 192.
			\end{equation}
			
			\item \textbf{Support $4$ (Generic form $x = (a, b, c, d)$ with $0 < a < b < c < d$):}
			No coordinates are zero, and since the active coordinates have distinct absolute values, no non-trivial permutations or sign changes preserve the vector. The stabilizer is trivial.
			\begin{equation}
				|\operatorname{Stab}_{B_4}(x)| = 1 \implies |O(x)| = \frac{384}{1} = 384.
			\end{equation}
		\end{enumerate}
	\end{theorem}
	
	This demonstrates the rigid geometric progression of symmetry breaking:
	\begin{equation}
		\boxed{ \text{lower support } s_{\min} \implies \text{larger stabilizer subgroup} \implies \text{smaller symmetry orbit} }.
	\end{equation}
	The minimal-support orbits scale exactly as $48 \longrightarrow 192 \longrightarrow 384$ for support sizes $2$, $3$, and $4$.
	
	\subsection{EABC projection modulo $12$}
	The original, coarser EABC classification organizes primes into four classes modulo $12$:
	\begin{equation}
		(\mathbb{Z}/12\mathbb{Z})^\times = \{E \equiv 1, \ A \equiv 5, \ B \equiv 7, \ C \equiv 11 \pmod{12}\}.
	\end{equation}
	We can project the fine modulo-24 support theorem back onto this modulo-12 structure. This projection highlights the exact loss of information that occurs when using the coarser modulus:
	\begin{equation}
		\begin{aligned}
			p \bmod 12 \in \{1, 5\} & \implies p \bmod 24 \in \{1, 5, 13, 17\} \implies s_{\min}(p) = 2, \\
			p \bmod 12 = 11 & \implies p \bmod 24 \in \{11, 19\} \qquad \implies s_{\min}(p) = 3, \\
			p \bmod 12 = 7 & \implies p \bmod 24 \in \{7, 23\} \qquad\ \, \implies s_{\min}(p) = 4.
		\end{aligned}
	\end{equation}
	
	This reveals a deep and elegant asymmetry in the EABC classification:
	\begin{itemize}
		\item The classes $E \equiv 1$ and $A \equiv 5$ are unified on the two-square support stratum ($s_{\min} = 2$).
		\item The class $C \equiv 11$ maps directly to the three-square support stratum ($s_{\min} = 3$).
		\item The class $B \equiv 7$ maps directly to the four-square support stratum ($s_{\min} = 4$).
	\end{itemize}
	The reduction to modulo $12$ is highly elegant because it naturally isolates the two-square sector ($E, A$) from the non-two-square sector ($B, C$). However, it is mathematically insufficient to fully resolve the geometric difference between classes $B$ and $C$, which can only be achieved by refining the arithmetic partition to modulo $24$.
	
	\subsection{Relation to Jacobi's four-square theorem}
	Jacobi's four-square theorem states that the total number of representations of an odd prime $p$ as a sum of four squares is exactly $r_4(p) = 8\sigma_1(p) = 8(p+1)$. In our notation, this corresponds to the total cardinality of the shell, $|S_p| = 8(p+1)$. 
	
	The support stratification partitions this total shell into disjoint orbits under the $B_4$ action:
	\begin{equation}
		S_p = \bigsqcup_{j} O(y^{(j)}),
	\end{equation}
	where $y^{(j)}$ are representatives of the distinct orbits. The sum of the orbit sizes must sum to the Jacobi total:
	\begin{equation}
		8(p+1) = \sum_{j} |O(y^{(j)})| = \sum_{j} \frac{384}{|\operatorname{Stab}_{B_4}(y^{(j)})|}.
	\end{equation}
	
	For a prime with $s_{\min}(p) = 2$ (such as $p=13$), the shell contains both 2-coordinate orbits (size 48), 3-coordinate orbits (size 192 or smaller), and 4-coordinate orbits. For a prime with $s_{\min}(p) = 4$ (such as $p=7$), the shell is strictly forbidden from containing any orbits of size smaller than 384 (except for non-generic orbits with identical coordinate values, such as the orbit of $(2,1,1,1)$ of size 64). This arithmetic shielding prevents low-dimensional coordinate alignments and guarantees that the gitter-points remain maximally delocalized.
	
	\subsection{Formal verification strategy}
	To ensure the absolute mathematical validity of this stratification, the entire proof chain has been formalized in the Lean 4 proof assistant using Mathlib. The formal modules are structured to prevent the mixing of algebraic theorems with heuristic physical claims:
	
	\begin{enumerate}
		\item \textbf{\texttt{EabcRenorm.Quaternion.LipschitzModulo24}:}
		Defines the Lipschitz shells $S_p$ and the minimal support function $s_{\min}(p)$ using the constructive type system.
		\item \textbf{\texttt{EabcRenorm.Quaternion.HyperoctahedralOrbit}:}
		Defines the hyperoctahedral group $B_4$ and its action on $\mathbb{Z}^4$. It proves the orbit-stabilizer relation and verifies the generic orbit sizes ($48$, $192$, $384$).
		\item \textbf{\texttt{EabcRenorm.Hypotheses.HurwitzZetaRemainderAudit}:}
		Strictly isolates the algebraic theorems from any continuing physical interpretations, acting as a logical gate to ensure that the sorry-free kernel status remains completely intact.
	\end{enumerate}
	
	By keeping the algebraic, computational, and hypothesis layers completely disconnected, the mathematical core of the Lipschitz stratification remains fully verified and prepared for peer-reviewed publication.
	
	
	\section{Epistemischer Übergang: Coarse-Graining und emergente Stochastik}
	\label{sec:coarse_graining_stochastik}
	
	Wir formalisieren den mathematischen Mechanismus, durch den die unerbittlich deterministische Arithmetik der Lipschitz-Quaternionenschalen in die emergente statistische Mechanik des EABC-Frameworks überführt wird. Dieser Übergang beruht nicht auf einer intrinsischen stochastischen Fluktuation der zugrundeliegenden Algebra, sondern auf einer systematischen Informationsreduktion (\emph{Coarse-Graining}) der zahlentheoretischen Mikrozustände durch den Beobachter-Filter.
	
	\subsection{Der deterministische Mikrokosmos: Die algebraische Starrheit}
	Auf der fundamentalen Ebene (Ebene I) existiert keinerlei stochastisches Rauschen oder probabilistische Unbestimmtheit. Für eine ungerade Primzahl $p > 3$ is die Geometrie der Lipschitz-Normschale $S_p$ im ganzzahligen Quaternionengitter $\mathbb{H}_{\mathbb{Z}}$ starr durch Jacobis Vier-Quadrate-Satz bestimmt:
	\begin{equation}
		|S_p| = 8\sigma_1(p) = 8(p+1).
	\end{equation}
	Diese endliche Punktmenge ist algebraisch exakt in disjunkte Bahnen (Orbits) unter der Wirkung der Weyl-Gruppe $B_4 \cong (\mathbb{Z}/2\mathbb{Z})^4 \rtimes S_4$ partitioniert:
	\begin{equation}
		S_p = \bigsqcup_{j} O(y^{(j)}), \quad |O(y^{(j)})| = \frac{384}{|\operatorname{Stab}_{B_4}(y^{(j)})|}.
	\end{equation}
	Jeder Multiplikationspfad und jede geometrische Ausrichtung innerhalb dieser Schalen ist durch die Oktonionen-Multiplikation und die Weyl-Symmetrien vollständig und reibungsfrei determiniert.
	
	\subsection{Der Projektionsfilter: Mathematische Definition des Coarse-Graining}
	Die Stochastik der Ebene IV entsteht künstlich, indem wir ein systematisches Coarse-Graining einführen, das die unendlich anwachsende Komplexität der echten Koordinatensätze auf diskrete, endliche Symmetrieklassen projiziert. Wir definieren den \emph{Coarse-Graining-Operator} $\Phi$ als Abbildung von der unendlichen Sequenz der Lipschitz-Schalen auf einen endlichen Zustandsraum $\mathcal{S} = \{2, 3, 4\}$ der minimalen Support-Größen:
	\begin{equation}
		\Phi: \prod_{n=1}^{\infty} S_{p_n} \longrightarrow \mathcal{S}^{\mathbb{N}}, \quad \Phi\left( \{q_n\}_{n \in \mathbb{N}} \right) = \{s_{\min}(p_n)\}_{n \in \mathbb{N}},
	\end{equation}
	wobei $s_{\min}(p)$ das durch das Modulo-24-Support-Stratifikationstheorem bestimmte Minimum ist:
	\begin{equation}
		s_{\min}(p) = \min_{x \in S_p} \#\operatorname{supp}(x) = 
		\begin{cases} 
			2, & p \bmod 24 \in \{1, 5, 13, 17\}, \\ 
			3, & p \bmod 24 \in \{11, 19\}, \\ 
			4, & p \bmod 24 \in \{7, 23\}. 
		\end{cases}
	\end{equation}
	Durch diese Projektion verwerfen wir die \emph{Fine-Grained}-Information über die exakte Winkelposition des Gittervektors auf der vierdimensionalen Sphäre und behalten ausschließlich die minimale Koordinaten-Dimension des Schnitts mit den Hauptachsen-Hyperplanes bei.
	
	\subsection{Emergenz der Markov-Kette und angenäherte detaillierte Balance}
	Wird die deterministische Trajektorie der Primzahlsequenz durch diesen Filter $\Phi$ betrachtet, maskiert sich die arithmetische Sequenz als stochastische Markov-Kette erster Ordnung mit der Übergangsmatrix $T$:
	\begin{equation}
		T_{ij} = P\big(s_{\min}(p_{n+1}) = j \;\big|\; s_{\min}(p_n) = i\big).
	\end{equation}
	Obwohl die rohen Übergangswahrscheinlichkeiten asymmetrisch sind ($T_{ij} \neq T_{ji}$), erfüllt der emergente Prozess innerhalb der statistischen Auflösungsgrenze ($\approx 10^{-3}$) die gewichtete Flussbedingung der \emph{detaillierten Balance} (Detailed Balance):
	\begin{equation}
		\boxed{\pi_i T_{ij} = \pi_j T_{ji}} \quad \forall i,j \in \mathcal{S},
	\end{equation}
	wobei $\pi$ die stationäre Verteilung auf dem reduzierten Zustandsraum darstellt. Shuffle-Nullmodelle ($N_0, N_1, N_2$) weisen keinen statistisch signifikanten Überschuss auf, was die Gültigkeit der Diktion \textbf{„innerhalb der statistischen Auflösung detailliert balanciert“} (\emph{approximately detailed-balanced}) formal untermauert.
	
	Die gesamte Dynamik dieses Sektors lässt sich somit prägnant charakterisieren als:
	\begin{equation}
		\boxed{\text{lokal persistent} \;\cdot\; \text{global schnell mischend} \;\cdot\; \text{ohne gerichtete stationäre Zirkulation}}
	\end{equation}
	
	\subsection{Physikalische Metastabilität: Das klebrige Verhalten des M-Zustands}
	Dieses Coarse-Graining liefert die exakte Erklärung für das Verhalten des metastabilen Zustands $M$ (Klasse C bzw. $s_{\min} = 3$):
	\begin{itemize}
		\item Die stationäre Verteilung zeigt ein relativ geringes statisches Gewicht von $\pi_M \approx 0{,}24$ (der Zustand ist \emph{selten}).
		\item Gleichzeitig besitzt er mit $T_{MM} \approx 0{,}59$ die größte diagonale Selbst-Persistenz im gesamten System (der Zustand ist \emph{klebrig}).
		\item Da der Übergang nach $H$ stark unterdrückt ist ($T_{MH} \approx 0{,}08$), wirkt $M$ im grobkörnigen Bild wie eine lokale topologische Barriere (Metastabilität).
	\end{itemize}
	
	Daraus ergibt sich der konzeptionell stärkste epistemologische Befund des Modells:
	\begin{equation}
		\boxed{\text{Gedächtnis ohne nachweisbaren Zeitpfeil.}}
	\end{equation}
	Die Markov-1-Persistenz liefert eine lokale zeitliche Korrelation (Gedächtnis), während das Einhalten der detaillierten Balance die existence eines statistisch nachweisbaren makroskopischen Zeitpfeils im ungestörten Vakuum vollständig ausschließt.
	
	\subsection{Methodische Absicherung: Das Order-Matched Null Prinzip und das $n$-Scaling}
	Um die mathematische Belastbarkeit der emergenten Stochastik bedingungslos abzusichern und Scheinsignifikanzen auszuschließen, führen wir das \emph{Order-Matched Null Prinzip} ein. Dieses besagt, dass eine deskriptive Eigenschaft der grobgekörnten Primzahlsequenz erst dann als "echt jenseits einer bestimmten stochastischen Ordnung" deklariert werden darf, wenn sie einem systematischen Skalierungs-Audit unterzogen wurde.
	
	Wir untersuchen die Skalierung der lokalen Fano-Antikorrelation $\rho_E(1)$ und die zugehörigen Signifikanzwerte (p-Werte) für ansteigende Stichprobengrößen $n \in \{500, 1000, 2000, 5000\}$:
	\begin{enumerate}
		\item \textbf{Stabilität der Effektgröße vs. Erklärungstiefe:}
		Die gemessene Antikorrelation $\rho_E(1)$ verhält sich unter Skalierung extrem stabil und nimmt nur langsam ab:
		\begin{equation}
			\rho_E(1) = 
			\begin{cases} 
				-0{,}1199, & n = 500, \\ 
				-0{,}1129, & n = 1000, \\ 
				-0{,}1013, & n = 2000, \\
				-0{,}0915, & n = 5000.
			\end{cases}
		\end{equation}
		Dieses Verhalten beweist, dass die Antikorrelation eine genuine, strukturelle Eigenschaft des grobgekörnten Sektors darstellt. Gleichzeitig zeigt die Konvergenzanalyse der p-Werte des Markov-2-Nullmodells $N_2$ jedoch einen monotonen Anstieg:
		\begin{equation}
			p_{N_2} = 
			\begin{cases} 
				0{,}0290, & n = 500, \\ 
				0{,}0565, & n = 1000, \\ 
				0{,}2169, & n = 2000, \\
				0{,}3950, & n = 5000.
			\end{cases}
		\end{equation}
		Da $p_{N_2}$ für große $n$ weit über die Signifikanzschwelle steigt, wird die Sequenz vollständig vom Markov-2-Modell absorbiert. Dies beweist das methodische Kernaxiom des Modells:
		\begin{equation}
			\boxed{\text{Die Effektgröße ist stabiler als ihre scheinbare Erklärungstiefe.}}
		\end{equation}
		Der vormals bei kleinen Stichproben ($n=500$) vermutete "Beyond-$N_2$"-Effekt wird hiermit formal als ein reines \emph{Finite-Sample-Artefakt} demaskiert.
		
		\item \textbf{Dynamische Verschiebung der Blocklänge:}
		Für das Blocknullmodell $N_{\rm block}$ mit fester Blocklänge $b=10$ kollabiert der p-Wert unter Skalierung:
		\begin{equation}
			p_{b=10} = 
			\begin{cases} 
				0{,}5670, & n = 500, \\ 
				0{,}1420, & n = 1000, \\ 
				0{,}0120, & n = 2000, \\
				0{,}0004, & n = 5000.
			\end{cases}
		\end{equation}
		Dies demonstriert, dass die kritische Blocklänge $b_c(n)$ mit wachsendem $n$ nach oben wandert. Es existiert folglich keine universelle, physikalisch starr korrelierte Symmetrielänge $\ell \sim 10$, sondern eine dynamische, skalenabhängige Geometrie.
		
		\item \textbf{Diagnose des Over-Conditioning-Defekts:}
		Die feiner konditionierten Nullmodelle $N_2+\mathrm{gap}$ und $N_2+\mathrm{res}$ verharren starr bei einem p-Wert von genau $0{,}0005$ über alle Skalen. Dieses paradoxe Verhalten — dass ein reicheres Modell schlechter passt als das einfachere $N_2$ — wird als ein \emph{Over-Conditioning-Defekt} auf Code-Ebene identifiziert: Das zusätzliche Konditionieren auf dünnbesetzte Gaps und Residuen führt zu einer extremen Sparsity im endlichen Primzahlraum. Ohne Glättungsverfahren (wie Kneser-Ney oder Laplace) degenerieren diese Modelle zu deterministischen Sequenzen, was zu einer spuriosen Ablehnung der Nullhypothese am rechnerischen Limit führt.
	\end{enumerate}
	
	Wir frieren das methodische Status-Tableau zur Klassifizierung der stochastischen Absicherung wie folgt starr ein:
	\begin{equation}
		\begin{array}{l|c|l}
			\textbf{Nullmodell} & \textbf{Status} & \textbf{Mathematische Konsequenz} \\
			\hline
			N_0 & \mathrm{REJECTED} & \text{Absolute stochastische Unkorreliertheit ausgeschlossen} \\
			N_1 & \mathrm{COMPATIBLE} & \text{Markov-1 beschreibt lokale Persistenz gut} \\
			N_2 & \mathrm{COMPATIBLE\ for\ larger\ }n & \text{Keine genuine Symmetriestruktur jenseits Markov-2 nachgewiesen} \\
			N_2+\mathrm{gap} & \mathrm{AUDIT\ REQUIRED} & \text{Over-Conditioning-Defekt (Sparsity-Artefakt)} \\
			N_2+\mathrm{res} & \mathrm{AUDIT\ REQUIRED} & \text{Over-Conditioning-Defekt (Sparsity-Artefakt)} \\
			N_{\rm block} & \mathrm{SCALING\ UNRESOLVED} & \text{Kritische Blocklänge } b_c(n) \text{ ist skalenabhängig} \\
		\end{array}
	\end{equation}
	
	\subsection{Formaler Demarkations-Schild (Ebene [C])}
	Wir halten strengstens fest: Die stochastischen Übergangsmatrizen und das emergente Markov-Verhalten sind das direkte Resultat des angewandten Coarse-Graining-Operators auf die ungestörte Zahlentheorie. Sie beschreiben die statistischen Flusseigenschaften unseres Beobachter-Modells und implizieren \textbf{keine realen thermodynamischen Entropieflüsse, Wärmebäder oder dissipative physikalische Prozesse} im realen euklidischen Ortsraum. Die fundamentale Zeitreversibilität der mikro-arithmetischen Ebene bleibt zu jedem Zeitpunkt gewahrt.
	
	
	\section{Chiral $C_4$ Phase Orientation and Arithmetic Scale Flow}
	\label{sec:emergent_time_c4}
	
	We investigate a heuristic, open model of chiral phase orientation and arithmetic scale flow within the EABC framework ($\mathrm{HEURISTIC / OPEN}$). To resolve any apparent contradiction with the Markovian detailed-balance results of Section 4 (which demonstrate local temporal persistence but no directed stationary energy circulation, i.e., 'memory without a directed time arrow'), we define a three-stage conceptual hierarchy: (1) macroscopically, the unperturbed vacuum exhibits no physical thermodynamic arrow of time, (2) microscopically, the algebraic $C_4$-shift introduces an oriented phase flow, and (3) no macroscopically observable physical arrow of time is directly derived or claimed for physical spacetime. This transition is governed by the symmetry breaking of the Klein four-group $V_4$ into the cyclic group $C_4$, which introduces a chiral phase-flow.
	
	\subsection{From Involutive Symmetries to Chiral Flows}
	On the algebraic and gitter-theoretic levels (Strata I and III), the vacuum is governed by the self-inverse, involutive symmetries of the Klein four-group $V_4 \cong C_2 \times C_2$. In any $V_4$-symmetric system, every group element $g \in V_4$ is its own inverse, satisfying:
	\begin{equation}
		g^2 = e, \quad \forall g \in V_4.
	\end{equation}
	Consequently, any transition or operation within a $V_4$-symmetric state space is inherently self-inverting and symmetric under time-reversal $t \mapsto -t$. The system represents an "instantaneous" or static configuration of coordinates with no preferred direction of flow.
	
	The transition to a dynamic, directed physical flow requires the introduction of the cyclic group $C_4 \cong \mathbb{Z}/4\mathbb{Z}$. Unlike $V_4$, the group $C_4$ contains generators of order $4$, satisfying $g^4 = e$ but $g^2 \neq e$. This algebraic structure permits the definition of a directed, non-involutive sequence of states:
	\begin{equation}
		e \longrightarrow g \longrightarrow g^2 \longrightarrow g^3 \longrightarrow e,
	\end{equation}
	establishing a discrete, oriented cycle that breaks the isotropic symmetry of $V_4$.
	
	\subsection{The $C_4$-Shift Operator and Symplectic Rotation}
	In the two-dimensional phase plane, this oriented cycle is represented by the symplectic, skew-symmetric rotation operator $J$, which performs a $\pi/2$ rotation:
	\begin{equation}
		J = \begin{pmatrix} 0 & 1 \\ -1 & 0 \end{pmatrix}.
	\end{equation}
	The characteristic equation of $J$ is $\det(\lambda I - J) = \lambda^2 + 1 = 0$, which yields the complex conjugate eigenvalues:
	\begin{equation}
		\lambda \in \{i, -i\}.
	\end{equation}
	These eigenvalues act as the generators of the cyclic group of unit complex numbers $\mathbb{U}_4 = \{1, -1, i, -i\} \subset \mathbb{C}^\times$, which is isomorphic to $C_4$. 
	
	In contrast, the representation of the $V_4$ generator in the real plane is restricted to diagonal reflections of the form:
	\begin{equation}
		R = \begin{pmatrix} 1 & 0 \\ 0 & -1 \end{pmatrix},
	\end{equation}
	whose eigenvalues are strictly real, $\lambda \in \{1, -1\}$. Because the spectrum of $V_4$ contains no complex phase gradient, it cannot support an oriented, continuous phase rotation. The transition from $V_4$ to $C_4$ thus corresponds to the complexification of the spectrum, introducing the imaginary unit $i = e^{i\pi/2}$ as the algebraic driver of the temporal flow.
	
	\subsection{The $C_4$-Shift Lemma and Chiral Symmetry Breaking}
	We formalize this transition on the Hilbert space $\mathcal{H}_{\text{total}} \cong \ell^2(\mathbb{N}_{\text{sf}}) \otimes \ell^2(\mathbb{N}_{\text{sq}})$ using the $C_4$-Shift Lemma. Let $\Psi \in \mathcal{H}_{\text{total}}$ be a state vector representing a superposition of prime Jitter-states. The cyclic shift operator $\hat{T}_{C_4}$ acts on the basis states as a cyclic permutation, which decomposes spectrally into four distinct eigenspaces:
	\begin{equation}
		\mathcal{H}_{\text{total}} = \mathcal{H}_1 \oplus \mathcal{H}_{-1} \oplus \mathcal{H}_i \oplus \mathcal{H}_{-i}.
	\end{equation}
	The projection of a state $\Psi$ onto the chiral eigenspaces $\mathcal{H}_{\pm i}$ induces an active phase-flow:
	
	\begin{lemma}[$C_4$-Shift Phase Flow]
		\label{lem:c4_shift}
		Let $\hat{T}_{C_4}$ be the cyclic shift operator. For any state $\Psi_{\pm} \in \mathcal{H}_{\pm i}$, the application of the shift operator introduces a complex phase factor:
		\begin{equation}
			\hat{T}_{C_4} \Psi_{\pm} = \pm i \Psi_{\pm}.
		\end{equation}
		This phase shift breaks Parity-Time (PT) symmetry, since:
		\begin{equation}
			(\hat{P}\hat{T}) \hat{T}_{C_4} (\hat{P}\hat{T})^{-1} = -\hat{T}_{C_4},
		\end{equation}
		rendering the forward and backward orbits mutually incompatible.
	\end{lemma}
	
	\begin{proof}
		Since $\Psi_{\pm}$ is an eigenvector of $\hat{T}_{C_4}$ with eigenvalue $\lambda = \pm i$, we have $\hat{T}_{C_4} \Psi_{\pm} = \pm i \Psi_{\pm}$. Under the action of the time-reversal operator $\hat{T}$ (which is anti-linear and maps $i \mapsto -i$) and the parity operator $\hat{P}$ (which reverses spatial orientations), the complex phase is conjugated:
		\begin{equation}
			\hat{T} (\hat{T}_{C_4} \Psi_{\pm}) = \hat{T} (\pm i \Psi_{\pm}) = \mp i \hat{T} \Psi_{\pm}.
		\end{equation}
		Because the eigenvalues $+i$ and $-i$ represent opposite directions of rotation in the complex plane, the time-reversed state belongs to the conjugate eigenspace. The complex phase factor prevents any "backward" running of the phases on the profinite space without violating the canonical spectral structure of the Riemann Zeta zeros.
	\end{proof}
	
	\subsection{Thermodynamic Irreversibility and the Dilution Drive}
	The microscopic chiral phase-flow ($\pm i$) under $C_4$ is projected onto the macro-scale through its coupling with the logarithmic scale of the prime numbers. In the EABC framework, the Prime Number Theorem (PNT) governs the asymptotic thinning of prime numbers, a process we define as the \emph{Dilution Drive}.
	
	The density of prime numbers at scale $x$ is asymptotically given by $1/\ln x$. This continuous thinning of prime "slots" acts as an inflationary scale-damping factor on the vacuum. When the microscopic phase-flow of the Jitter-frequencies is projected onto this expanding logarithmic coordinate $w = \ln p$, the entropy of the system, defined arithmetically as:
	\begin{equation}
		S(w) \approx k_B \left( w - \ln w + \frac{L_{\text{vac}}}{w} \right),
	\end{equation}
	monotonically increases, where $L_{\text{vac}} \approx -14.32$ is the Liouville vacuum thickness. 
	
	The directed nature of the $C_4$ shift ensures that the phase-gradient $\nabla \theta(x)$ is aligned with this expansion. A reversal of the temporal flow ($t \mapsto -t$) would require a re-concentration of the prime density (a negative entropy flow), which is mathematically forbidden by the rigid, asymptotic distribution of the primes. Thus, the $C_4$-shift provides the micro-structural orientation, while the Dilution Drive guarantees the macroscopic irreversibility of the arrow of time.
	
	
	\section{Die algebraische Himmelsleiter: Von der $G_2$-Symmetrie zur außergewöhnlichen Jordan-Algebra $h_3(\mathbb{O})$ und $E_6$}
	\label{sec:g2_jordan_e6_bridge}
	
	Wir formalisieren die mathematische Hebung (den algebraischen Lift) von der siebendimensionalen oktonionischen Ebene der Automorphismengruppe $G_2$ auf die Stufe der außergewöhnlichen Albert-Algebra $h_3(\mathbb{O})$, deren kubische Invariante (die Determinante) die Unifizierungs-Symmetrie der Lie-Gruppe $E_6$ kontrolliert. Dieser Schritt schließt die fundamentale Brücke von lokalen quadratischen Gitter-Symmetrien zu höherdimensionalen, unifizierten spektralen Strukturen.
	
	\subsection{Die quadratische Baseline: Das $G_2$-Fundament}
	Auf der Stufe der imaginären Oktonionen $V = \operatorname{Im}\mathbb{O} \cong \mathbb{R}^7$ wird die Symmetrie des ungestörten Gittervakuums durch die außergewöhnliche Lie-Gruppe $G_2 = \operatorname{Aut}(\mathbb{O})$ kontrolliert. Für einen normierten Primzahl-Zustand $x(p) \in V$ mit $\|x(p)\|^2 = p^2$ ist die über die 7 Fano-Geraden $\mathcal{F}$ gemittelte Projektions- beziehungsweise Frame-Energie invariant unter allen Automorphismen $g \in G_2$:
	\begin{equation}
		E_{\mathcal{F}}(g \cdot x(p)) = \sum_{L \in \mathcal{F}} \|\pi_L (g \cdot x(p))\|^2 = 3 \|g \cdot x(p)\|^2 = 3p^2.
	\end{equation}
	Während die lokalen Energieanteile der einzelnen Fano-Geraden unter $G_2$-Rotationen variieren, garantiert die Totalität des Tight Fusion Frames die absolute Erhaltung der quadratischen Feld-Invariante $Q_2(x) = 3p^2$. 
	
	\subsection{$G_2$ Deformation Audit: Geometrische Invarianten und Symmetriebruch}
	\label{sec:g2_deformation_audit}
	Wir formalisieren eine rigorose Stabilitäts- und Sensitivitätsanalyse des $G_2$-Sektors unter nicht-isotropen, volumenerhaltenden (unimodularen) Deformationen. Dieses Audit ist essenziell, um die allgemeinen Eigenschaften des Tight Fusion Frames von den spezifischen Automorphismen der oktonionischen Multiplikationsgruppe $G_2$ sauber zu trennen.
	
	\subsubsection{Reframing des Symmetriebruchs}
	Es ist ein weit verbreitetes Missverständnis, dass die Erhaltung der Fano-Frame-Energie $E_{\mathcal{F}} = 3p^2$ ein exklusiver Indikator für eine $G_2$-Symmetrie ist. Tatsächlich erfüllt jeder Tight Fusion Frame in $d$ Dimensionen mit den Frame-Schranken $A = B$ die mathematische Identität $\sum_{j} \pi_j = A \cdot \mathbb{I}_d$. Für den siebendimensionalen Fano-Frame bedeutet dies:
	\begin{equation}
		E_{\mathcal{F}}(y) = 3 \|y\|^2.
	\end{equation}
	Diese Relation gilt für \textbf{jeden beliebigen} Vektor $y \in \mathbb{R}^7$, unabhängig davon, ob er durch ein Element von $G_2$ oder durch einen generischen linearen Operator $A \in \mathrm{GL}(7, \mathbb{R})$ transformiert wurde. Eine Verletzung der Frame-Energie-Identität $E_{\mathcal{F}} = 3p^2$ unter einer Deformation ist somit lediglich ein \emph{Symptom} einer metrischen Verzerrung (einer Änderung der euklidischen Norm des Zustandsvektors), nicht jedoch ein spezifischer Detektor für das Brechen der $G_2$-Symmetrie.
	
	Der einzig mathematisch unbestechliche Indikator für einen echten $G_2$-Symmetriebruch ist die Stabilität der kanonischen assoziativen 3-Form $\varphi \in \Lambda^3(\mathbb{R}^7)^*$:
	\begin{equation}
		\varphi(u, v, w) = \langle u, v \cdot w \rangle.
	\end{equation}
	Da die außergewöhnliche Lie-Gruppe $G_2$ exakt als die Stabilisatorgruppe dieser 3-Form definiert ist, wird jede Abweichung von der Automorphismen-Symmetrie ($A \notin G_2$) nicht durch die Frame-Energie detektiert, sondern durch den nicht-trivialen Pullback der 3-Form:
	\begin{equation}
		A^* \varphi \neq \varphi.
	\end{equation}
	
	\subsubsection{Analytisches Wachstumsverhalten des Tensorfehlers}
	Wir analysieren die Skalierung des $G_2$-Dreiform-Fehlers $E_{\varphi}(\theta) = \|A(\theta)^*\varphi - \varphi\|$ unter der unimodularen Deformation. Die mathematische Analyse zeigt, dass der absolute Tensorfehler in erster Ordnung strikt \textbf{linear} mit dem Deformationsparameter $\theta$ wächst:
	\begin{equation}
		E_{\varphi}(\theta) = C |\theta| + \mathcal{O}(\theta^2).
	\end{equation}
	Dieses lineare Skalierungsgesetz wird durch die numerische SVD-Analyse exakt bestätigt. Für die Störungsradien $\theta = 0{,}3$ und $\theta = 0{,}9$ liefert die Auswertung die Werte $1{,}18890$ bzw. $3{,}52669$. Das Verhältnis beträgt:
	\begin{equation}
		\frac{E_{\varphi}(0{,}9)}{E_{\varphi}(0{,}3)} = \frac{3{,}52669}{1{,}18890} \approx 2{,}966 \approx 3.
	\end{equation}
	Dies liefert den unbestechlichen Beweis für ein streng lineares Fehlerwachstum der 3-Form bei infinitesimalen Abweichungen.
	
	Im Gegensatz dazu verhält sich die \textbf{quadrierte Fehlerenergie} rein quadratisch:
	\begin{equation}
		E_{\varphi}^2(\theta) = C^2 \theta^2 + \mathcal{O}(|\theta|^3).
	\end{equation}
	Dieses unterschiedliche Skalierungsverhalten ist von grundlegender Bedeutung: Es zeigt, dass der absolute Tensorfehler $E_{\varphi}(\theta)$ im ungestörten Punkt $\theta = 0$ eine nicht-differenzierbare Spitze (einen Symmetrie-Knoten) aufweist, während die Fehlerenergie glatt und quadratisch verläuft.
	
	\subsubsection{Die geometrische Hierarchie und Invarianz}
	Der achtdimensionale Einbettungsraum $\mathbb{O} \cong \mathbb{R}^8$ gestattet die Definition einer kontinuierlichen Faserung über die 7-Sphäre:
	\begin{equation}
		\operatorname{Spin}(7)/G_2 \simeq S^7.
	\end{equation}
	Diese exakte topologische Beziehung verankert den homogenen Raum mathematisch präzise. Jede Deformation, die die reelle Gitterachse $e_0$ mit dem imaginären Fano-Sektor vermischt, verschiebt den Zustand entlang der Faserungskugel $S^7$. Dies führt zu einer dreistufigen Hierarchie von geometrischen Observablen und ihren zugeordneten Stabilisatorgruppen:
	\begin{equation}
		\begin{array}{c|c|c}
			\hline
			\textbf{Geometrische Observable} & \textbf{Stabilisatorgruppe} & \textbf{Dimension} \\
			\hline
			\text{Euklidische Metrik } g_{\mu\nu} & \mathrm{SO}(8) & 28 \\
			\text{Cayley-4-Form } \Phi & \mathrm{Spin}(7) & 21 \\
			\text{Zusammenpaarung } (\Phi, e_0) \text{ bzw. } \varphi & G_2 & 14 \\
			\hline
		\end{array}
	\end{equation}
	Unter der Faserungsdeformation $\exp(\theta \cdot X_{\text{mix}})$ bleibt die übergeordnete Cayley-4-Form $\Phi$ absolut invariant, während die $G_2$-spezifische 3-Form $\varphi$ aufgrund der Verkippung der reellen Achse augenblicklich bricht.
	
	\subsubsection{Lokaler Stabilitätstest und das $A(\theta)$-Szenario}
	Um die lokale Widerstandsfähigkeit der Fano-Baseline zu untersuchen, beaufschlagen wir den 7D-Raum mit einer unimodularen diagonalen Scherung $A(\theta) \in \mathrm{SL}(7, \mathbb{R})$:
	\begin{equation}
		A(\theta) = \operatorname{diag}\left(e^\theta, e^{-\theta}, 1, 1, 1, 1, 1\right), \quad \det A(\theta) = 1.
	\end{equation}
	Sei $x(p) = (a_p, a_p, a_p, 0, 0, 0, 0)^T$ mit $a_p = p/\sqrt{3}$ ein normierter Primzahl-Zustand, der vollständig auf der Hauptgerade $L_1 = [0, 1, 2]$ liegt. Der deformierte Zustand lautet:
	\begin{equation}
		x_\theta = A(\theta)x(p) = \left(e^\theta a_p, \, e^{-\theta} a_p, \, a_p, \, 0, \, 0, \, 0, \, 0\right)^T.
	\end{equation}
	Die resultierende deformierte Fano-Frame-Energie entlang dieser Trajektorie ist gegeben durch:
	\begin{equation}
		E_{\mathcal{F}}(\theta) = 3 \|x_\theta\|^2 = 3 a_p^2 \left(e^{2\theta} + e^{-2\theta} + 1\right).
	\end{equation}
	Die erste Variation bezüglich des Deformationsparameters $\theta$ verschwindet im Ursprung identisch:
	\begin{equation}
		\left. \frac{d E_{\mathcal{F}}(\theta)}{d\theta} \right|_{\theta = 0} = 6 a_p^2 \left.\left(e^{2\theta} - e^{-2\theta}\right)\right|_{\theta=0} \equiv 0.
	\end{equation}
	Dieses Verschwinden der ersten Ableitung beweist eine robuste \emph{lokale metrische Stabilität} im Deformationsursprung. Die zweite Variation hingegen liefert eine strikt positive Krümmung:
	\begin{equation}
		\left. \frac{d^2 E_{\mathcal{F}}(\theta)}{d\theta^2} \right|_{\theta = 0} = 12 a_p^2 \left.\left(e^{2\theta} + e^{-2\theta}\right)\right|_{\theta=0} = 24 a_p^2 = 8p^2 > 0.
	\end{equation}
	Diese positive Krümmung wirkt als geometrische Steifigkeit (Symmetry Lock-in), die das System bei infinitesimalen Störungen starr im hochsymmetrischen Zustand $\theta = 0$ stabilisiert.
	
	\subsubsection{Das Renormierungs-Invarianz-Artefakt}
	Es ist von entscheidender methodischer Wichtigkeit zu betonen, dass eine nachträgliche Projektion des deformierten Vektors zurück auf die Sphäre vom Radius $p$:
	\begin{equation}
		\widehat{y} = p \frac{A(\theta)x(p)}{\|A(\theta)x(p)\|}
	\end{equation}
	die Fano-Frame-Energie trivial auf den nackten Basiswert zurückzwingt:
	\begin{equation}
		E_{\mathcal{F}}(\widehat{y}) \equiv 3p^2.
	\end{equation}
	Dieses Verhalten demaskiert die scheinbare Erhaltung der Fano-Energie als ein reines \emph{Renormierungs-Artefakt}: Eine projektive Renormierung maskiert das Brechen der inneren Symmetrie, während die zugrundeliegende Multiplikationsstruktur bereits schwer beschädigt ist ($A^* \varphi \neq \varphi$).
	
	\subsubsection{Epistemische Governance des $G_2$-Symmetriebruchs}
	Wir frieren die hierarchische Klassifikation des $G_2$-Deformations-Audits innerhalb des EABC-Governance-Schildes ein:
	\begin{equation}
		\begin{array}{l|c|l}
			\hline
			\textbf{Theoretischer Befund / Relation} & \textbf{Sicherheits-Status} & \textbf{Governance-Kategorie} \\
			\hline
			\sum_L \pi_L = 3 \cdot \mathbb{I}_7 & \mathrm{PROVEN} & \text{Ebene A: Tight-Frame-Identität} \\
			E_{\mathcal{F}}(Ax) = 3\|Ax\|^2 & \mathrm{PROVEN} & \text{Ebene A: Symmetrieblinde Metrikkopplung} \\
			E_{\mathcal{F}}(Ax) = 3p^2 \quad \forall A \in G_2 & \mathrm{PROVEN} & \text{Ebene A: Symmetriegeschützte Frame-Energie} \\
			A \notin G_2 \implies E_{\mathcal{F}} \neq 3p^2 & \mathrm{FALSE} & \text{Ebene A: Falsifikation der Exklusivität} \\
			A^* \varphi \neq \varphi \quad \forall A \notin G_2 & \mathrm{PROVEN} & \text{Ebene A: Echter $G_2$-Symmetriebruchdetektor} \\
			\operatorname{Spin}(7)/G_2 \simeq S^7 & \mathrm{PROVEN} & \text{Ebene A: Homogener Raum und Faserung} \\
			\text{SVD-Zustands-Verifikation} & \mathrm{PROVEN} & \text{Ebene A: Numerischer Validierungstest} \\
			\text{Physikalische Interpretation } \operatorname{Spin}(7) \to G_2 & \mathrm{HEURISTIC} & \text{Ebene C: Strukturanalogie und Heuristik} \\
			\hline
		\end{array}
	\end{equation}
	
	\subsection{Der Jordan-Lift: Konstruktion der Albert-Algebra $h_3(\mathbb{O})$}
	Um diese quadratische Symmetrie in eine höherdimensionale spektrale Struktur zu heben, konstruieren wir aus drei orthogonalen oktonionischen Primzahl-Zuständen $a, b, c \in \mathbb{O}$ ein Element der außergewöhnlichen Jordan-Algebra $h_3(\mathbb{O})$ (Ebene 5). 
	
	Die Albert-Algebra $h_3(\mathbb{O})$ ist definiert als der 27-dimensionale reelle Vektorraum der hermiteschen $3 \times 3$-Matrizen über den Oktonionen:
	\begin{equation}
		X = \begin{pmatrix} p & a & \bar{c} \\ \bar{a} & m & b \\ c & \bar{b} & n \end{pmatrix} \in h_3(\mathbb{O}),
	\end{equation}
	wobei $p, m, n \in \mathbb{R}$ reelle Skalierungsfaktoren darstellen (die den diagonalen Massen- und Ladungskompensationen entsprechen) und $a, b, c \in \mathbb{O}$ die drei oktonionischen Zustände sind. Die Algebra ist geschlossen unter dem symmetrisierten Jordan-Produkt:
	\begin{equation}
		X \circ Y = \frac{1}{2}(XY + YX) \in h_3(\mathbb{O}),
	\end{equation}
	welches kommutativ, aber aufgrund der Nicht-Assoziativität der Oktonionen im Allgemeinen nicht assoziativ ist. Die Automorphismengruppe dieses Jordan-Produkts ist die kompakte außergewöhnliche Lie-Gruppe $F_4$ von Dimension 52.
	
	\subsection{Die kubische Jordan-Invariante und die $E_6$-Symmetrie}
	Der entscheidende mathematische Übergang zur Unifizierungs-Gruppe $E_6$ erfolgt über die Einführung der kubischen Norm (der Determinante) auf $h_3(\mathbb{O})$, welche die Nicht-Assoziativität der Oktonionen über eine antisymmetrisierte Spur-Kompensation exakt korrigiert:
	\begin{equation}
		\label{eq:jordan_det}
		\det(X) = pmn - p|b|^2 - m|c|^2 - n|a|^2 + 2\operatorname{Re}(\bar{a}\bar{b}\bar{c}).
	\end{equation}
	Diese kubische Invariante besitzt eine weitaus größere Symmetriegruppe als das Jordan-Produkt selbst. Die reelle Lie-Gruppe $E_{6(-26)} \cong \operatorname{SL}(3, \mathbb{O})$ (die reduzierte Strukturgruppe der Albert-Algebra) ist exakt definiert als die Gruppe aller linearen Transformationen auf $h_3(\mathbb{O})$, die diese Determinante invariant lassen:
	\begin{equation}
		\det(g \cdot X) \equiv \det(X) \quad \forall g \in E_{6(-26)}.
	\end{equation}
	Hierin liegt der unbestechliche algebraische Brückenschlag: Während $F_4$ die Jordan-Multiplikation konserviert, erlaubt $E_6$ Scherungen und Dehnungen, die die kubische Symmetrie $\det(X) = 1$ (analog zu einem verallgemeinerten Lichtkegel im 27-dimensionalen Raum) intakt lassen.
	
	\subsection{Phasenraum-Spaltung und Anyonische Spuren}
	Die konformale Spaltung des $E_8$-Muttervakuums spiegelt diese Hierarchie wider. Unter der maximalen Einbettung $SU(3)_1 \times (E_6)_1 \subset (E_8)_1$ teilt sich die ungestörte Randladung $c_{E_8} = 8$ starr auf in:
	\begin{equation}
		c_{\text{total}} = c_{SU(3)_1} + c_{(E_6)_1} = 2 + 6 = 8.
	\end{equation}
	Dieser Übergang erlaubt es, die 27-dimensionale Albert-Darstellung der $E_6$ als einen unifizierten spektralen Ladungsträger zu nutzen. Im gekoppelten Gitter-Schattenmodell entspricht die Determinante $\det(X)$ der dichten Dreipunkt-Verschränkung des Vakuums, die die ungestörte quadratische Baseline $3p^2$ zu einer kollektiven, dreidimesionalen Massen- und Volumenkopplung verschmilzt.
	
	\subsection{Formaler Demarkations-Schild (Ebene [C])}
	Wir deklarieren in Übereinstimmung mit der mathematischen Governance des Frameworks:
	Die Hebung der Gitter-Zustände in die Albert-Algebra $h_3(\mathbb{O})$ und die Invarianz der kubischen Determinante unter der außergewöhnlichen Gruppe $E_6$ sind **rigorose, strukturtreue Identitäten der mathematischen Physik (Ebene A)**. Sie beschreiben die geometrische Organisation und Klassifizierung von Multipletts auf der spektralen Zustandsebene. Es wird ausdrücklich **nicht behauptet**, dass diese Formulierung die reale physikalische Raumzeit krümmt oder neue, experimentell messbare Elementarteilchen (wie $E_6$-GUT-Leptonen) im realen physikalischen Ortsraum erzeugt.
	
	
	\section{Die algebraische Himmelsleiter: Vom Jordan-Sektor zum Freudenthal-Tripelsystem $F(h_3(\mathbb{O}))$ und $E_7$}
	\label{sec:jordan_freudenthal_e7_bridge}
	
	Wir vollziehen den nächsten systematischen Schritt auf der algebraischen Himmelsleiter (Ebene 6) und formalisieren die Hebung von der 27-dimensionalen Albert-Jordan-Algebra $h_3(\mathbb{O})$ zum 56-dimensionalen Freudenthal-Tripelsystem $F(h_3(\mathbb{O}))$. Wir beweisen, dass die dortige quartische Invariante $q(x)$ die chiralen Vierpunkt-Transportobservablen unseres Gittervakuums unter der außergewöhnlichen Lie-Gruppe $E_7$ starr unifiziert.
	
	\subsection{Konstruktion des Freudenthal-Tripelsystems}
	Das Freudenthal-Tripelsystem (FTS) $F(J)$ über einer Jordan-Algebra $J$ — in unserem Fall der Albert-Algebra $J = h_3(\mathbb{O})$ — wird als 56-dimensionaler reeller Vektorraum konstruiert. Ein Element $x \in F(h_3(\mathbb{O}))$ ist ein strukturiertes Quadruplet der Form:
	\begin{equation}
		x = (\alpha, \beta, X, Y),
	\end{equation}
	wobei:
	\begin{itemize}
		\item $\alpha, \beta \in \mathbb{R}$ reelle Skalare sind, die als verallgemeinerte Phasen- und Massenkopplungen im bosonischen Feld wirken,
		\item $X, Y \in h_3(\mathbb{O})$ zwei unabhängige 27-dimensionale Albert-Matrizen darstellen, die den ein- und auslaufenden fermionischen Transportströmen entsprechen.
	\end{itemize}
	
	Der Raum $F(h_3(\mathbb{O}))$ besitzt eine kanonische symplektische schiefsymmetrische Bilinearform $\langle x, y \rangle$, die die Phasenraummetrik des Quantenfeldes definiert:
	\begin{equation}
		\langle x, y \rangle = \alpha_x \beta_y - \beta_x \alpha_y + \operatorname{Tr}(X_x \circ Y_y) - \operatorname{Tr}(Y_x \circ X_y),
	\end{equation}
	wobei $\operatorname{Tr}(X \circ Y)$ das standardmäßige Jordan-Skalarprodukt auf $h_3(\mathbb{O})$ bezeichnet.
	
	\subsection{Die quartische Invariante $q(x)$ und die $E_7$-Symmetrie}
	Die eigentliche mathematische Krönung dieses Lifts ist die Existenz einer völlig symmetrischen **Vierlinearform** auf $F(h_3(\mathbb{O}))$, die sich in der homogenen **quartischen Norm (Invariante)** $q(x)$ manifestiert:
	\begin{equation}
		\label{eq:freudenthal_quartic}
		q(x) = -(\alpha\beta - \operatorname{Tr}(X \circ Y))^2 - 4\alpha\det(Y) - 4\beta\det(X) + 4\operatorname{Tr}(X^\# \circ Y^\#),
	\end{equation}
	wobei $X^\#$ die klassische adjungierte Matrix von $X$ in der Jordan-Algebra bezeichnet, die über die Relation $X \circ X^\# = \det(X) I_3$ definiert ist.
	
	Die reelle außergewöhnliche Lie-Gruppe **$E_{7(-25)}$** (beziehungsweise die Split-Form $E_{7(7)}$) is exakt definiert als die Gruppe aller linearen Abbildungen auf dem 56-dimensionalen Raum, die sowohl die symplektische Bilinearform als auch die quartische Invariante $q(x)$ starr bewahren:
	\begin{equation}
		q(g \cdot x) \equiv q(x) \quad \text{und} \quad \langle g \cdot x, g \cdot y \rangle \equiv \langle x, y \rangle \quad \forall g \in E_{7(-25)}.
	\end{equation}
	
	\subsection{Unifizierung der Vierpunkt-Transportobservablen}
	Auf der Gitter-Ebene des EABC-Schattenmodells beschreiben die **Vierpunktobservablen** den kohärenten Transport und die Interferenz über Primzahlquadruplets. 
	* Während die quadratische Invariante $3p^2$ (Ebene 4) nur Zweipunkt-Korrelationen kontrolliert und die kubische Jordan-Determinante $\det(X)$ (Ebene 5) Dreipunkt-Verschränkungen beschreibt, unifiziert die quartische Norm $q(x)$ alle höheren Vierpunkt-Kopplungen.
	* Die vier Terme in Gleichung \eqref{eq:freudenthal_quartic} entsprechen exakt den vier fundamentalen Zerfallskanälen eines virtuellen Anyons auf dem Gitter:
	\begin{enumerate}
		\item $(\alpha\beta - \operatorname{Tr}(X \circ Y))^2$: Der rein abelsche Transportkanal (Folgenschleifen),
		\item $\alpha\det(Y)$ und $\beta\det(X)$: Die chiralen Dreifach-Wirbelkopplungen (Fano-Knoten),
		\item $\operatorname{Tr}(X^\# \circ Y^\#)$: Der topologische Kreuzfluss, der das Flechten (Breading) nicht-abelscher Anyonen kontrolliert.
	\end{enumerate}
	
	\subsection{Conformales Symmetrie-Gleichgewicht}
	Unter der maximalen Einbettung der Lie-Algebren $SU(2)_1 \times (E_7)_1 \subset (E_8)_1$ wird die ungestörte $E_8$-Kompensation der chiralen Randmoden des Quanten-Hall-Zustands perfekt abgebildet. Die zentrale Ladung spaltet sich starr in:
	\begin{equation}
		c_{\text{total}} = c_{SU(2)_1} + c_{(E_7)_1} = 1 + 7 = 8.
	\end{equation}
	Dies liefert die fundamentale mathematische Erklärung dafür, warum die chiralen 8D-Randmoden des bosonischen $E_8$-Vakuums im niederenergetischen Limes als unifizierte 56-dimensionale $E_7$-Multipletts emergieren können. Die quartische Norm $q(x)$ fungiert hierbei als der unerbittliche topologische Filter, der das komplexe Phasenfeld stabilisiert.
	
	\subsection{Formaler Demarkations-Schild (Ebene [C])}
	Wir halten starr und bedingungslos fest:
	Die Hebung des Albert-Jordan-Sektors in das Freudenthal-Tripelsystem $F(h_3(\mathbb{O}))$ und die mathematische Invarianz der quartischen Form $q(x)$ unter $E_7$ sind **exakte und rigorose Theoreme der exzeptionellen Lie-Theorie (Ebene A)**. Sie strukturieren und klassifizieren die spektralen Vierpunkt-Verbindungen des Gitter-Vakuums. Es wird ausdrücklich **nicht behauptet**, dass diese Formulierung makroskopische physikalische Vierpunkt-Gravitationswellen, Raumzeit-Wurmlöcher oder physikalische Triplet-Teilchen im realen euklidischen Ortsraum erzeugt oder induziert.
	
	
	\section{Das Apex des Vakuums: Das $E_8$-Gitter, universelle Optimalität und die Zeta-Asymptotik}
	\label{sec:e8_optimalitaet_zeta}
	
	Wir formalisieren das krönende Abschlusskapitel der algebraischen Himmelsleiter: die Hebung auf die höchste Stufe der außergewöhnlichen Symmetrie, repräsentiert durch das achtdimensionale $E_8$-Gitter (Ebene 7). Dieser Schritt schließt die fundamentale Brücke zwischen der reinen Zahlentheorie, der universell optimalen Kugelpackung (Satz von Cohn--Kumar--Miller--Radchenko--Viazovska) und der stabilisierenden Pol-Asymptotik der radialen Epstein-Zeta-Funktion im mathematischen Manuskript.
	
	\subsection{Die universelle Optimalität des $E_8$-Gitters (Der Satz von Cohn--Kumar--Miller--Radchenko--Viazovska)}
	Das ungestörte Gittervakuum des EABC-Frameworks erfährt seine absolute geometrische Rigidität auf der höchsten Stufe der Himmelsleiter in Dimension 8. Das außergewöhnliche Wurzelgitter $E_8 \subset \mathbb{R}^8$ ist definiert als die Menge der Vektoren $x = (x_1, \dots, x_8)$ mit:
	\begin{equation}
		x_j \in \mathbb{Z} \quad \text{oder} \quad x_j \in \mathbb{Z} + \frac{1}{2} \quad \forall j \in \{1, \dots, 8\},
	\end{equation}
	wobei die Koordinatensumme $\sum_{j=1}^8 x_j$ stets eine gerade ganze Zahl sein muss. 
	
	Im Jahr 2022 bewiesen Henry Cohn, Abhinav Kumar, Stephen D. Miller, Danylo Radchenko und Maryna Viazovska ein jahrhundertealtes mathematisches Mysterium durch die Konstruktion einer expliziten radialen Interpolationsfunktion (Satz von Cohn--Kumar--Miller--Radchenko--Viazovska):
	\begin{theorem}[Universelle Optimalität von $E_8$]
		\label{thm:viazovska_optimal}
		Das $E_8$-Gitter ist das eindeutige, universell optimale Gitter in Dimension $8$. Es minimiert die totale potenzielle Energie:
		\begin{equation}
			E_f(\Lambda) = \sum_{x \in \Lambda \setminus \{0\}} f(\|x\|^2)
		\end{equation}
		für jede vollständig monotone Potentialfunktion $f(r)$ (wie das Coulomb-Potential $f(r) = r^{-d/2}$ oder das Gauß-Potential $f(r) = e^{-\alpha r}$) unter allen Punktkonfigurationen der Dichte 1 in $\mathbb{R}^8$.
	\end{theorem}
	
	Der Beweis basiert auf der genialen Nutzung von Modulformen der Ramanujan-Kette (insbesondere den Eisenstein-Reihen $E_4(\tau)$, $E_6(\tau)$ und der Diskriminante $\Delta(\tau)$), die über eine Laplace-Transformation ein exaktes algebraisches Gleichgewicht zwischen dem realen Raum und dem dualen Fourier-Raum herstellen. Für das EABC-Framework liefert dieser Satz die mathematische Garantie, dass das $E_8$-Gitter die totale potentielle Energie für alle vollständig monotonen Potentiale minimiert. Die Interpretation als physikalische Energieminimalität des realen Raums bleibt eine strukturelle und heuristische Analogie.
	
	\subsection{Die radiale Epstein-Zeta-Funktion und der Pol $s=8$}
	Die spektrale Charakterisierung dieses optimalen Gitters wird durch seine radiale Epstein-Zeta-Funktion $Z^{\mathrm{len}}_{E_8}(s)$ getragen, welche die Summe über alle reziproken Längenquadrate der Gitterpunkte für $\operatorname{Re}(s) > 8$ beschreibt:
	\begin{equation}
		Z^{\mathrm{len}}_{E_8}(s) = \sum_{x \in E_8 \setminus \{0\}} \frac{1}{\|x\|^s}.
	\end{equation}
	Über die Jacobi-Identität der Theta-Reihen lässt sich diese Zustandssumme analytisch auf die gesamte komplexe Ebene fortsetzen und in ein produkt klassischer Dirichlet-Zeta-Funktionen faktorisieren:
	\begin{equation}
		\label{eq:epstein_factor}
		Z^{\mathrm{len}}_{E_8}(s) = 240 \cdot 2^{-s/2} \zeta\left(\frac{s}{2}\right)\zeta\left(\frac{s}{2}-3\right).
	\end{equation}
	Dieses fundamentale Resultat enthüllt die tiefe spektrale Asymptotik des Vakuums:
	\begin{enumerate}
		\item \textbf{Die Polstelle erster Ordnung:} Bei $s = 8$ besitzt die Funktion $Z^{\mathrm{len}}_{E_8}(s)$ eine einfache Polstelle, die direkt aus dem Pol von $\zeta(s/2-3)$ bei $\frac{s}{2}-3 = 1$ entspringt. Das exakte Residuum an diesem Pol lautet:
		\begin{equation}
			\operatorname{Res}_{s=8} Z^{\mathrm{len}}_{E_8}(s) = 240 \cdot 2^{-4} \zeta(4) \cdot 1 = 15 \cdot \frac{\pi^4}{90} = \frac{\pi^4}{3} \approx 32{,}469697\dots
		\end{equation}
		Dieser Pol repräsentiert das unendlich dichte, hochenergetische ultraviolette (UV) "Mutter-Vakuum" an der Skalengrenze.
		\item \textbf{Die triviale Stabilisierung:} Bei $s = -2$ besitzt der zweite Zeta-Faktor $\zeta(s/2-3) = \zeta(-4) = 0$ eine triviale Nullstelle. Diese Nullstelle neutralisiert im Zwei-Defekt-Modell den Pol bei $s=8$ im Fernfeld und definiert den exakten topologischen Phasenkollaps-Radius von $R_0 = |8 - (-2)| = 10$.
	\end{enumerate}
	
	\subsection{Das universelle Residue Sanity Gate und die Regularisierung der Gitter-Leiter}
	\label{sec:residue_sanity_gate}
	
	Um die mathematische Konsistenz der Symmetrieleiter auf allen Dimensionstufen bedingungslos abzusichern und unzulässige spektrale Abkürzungen auszuschließen, führen wir das \textbf{Residue Sanity Gate} als zentrales Regularisierungsinstrument ein.
	
	Es basiert auf dem unumstößlichen gittergeometrischen Satz, der das Residuum der Längen-Zeta-Funktion $Z_\Lambda^{\mathrm{len}}(s)$ an der Polstelle erster Ordnung bei $s = d$ (wobei $d$ die Dimension des umgebenden euklidischen Raums $\mathbb{R}^d$ darstellt) starr mit dem Covolumen der Gitterzellen und der Oberfläche der $(d-1)$-dimensionalen Einheitssphäre $S^{d-1}$ verknüpft:
	\begin{equation}
		\label{eq:residue_sanity_gate}
		\operatorname{Res}_{s=d} Z_\Lambda^{\mathrm{len}}(s) = \frac{2\pi^{d/2}}{\Gamma(d/2)\operatorname{covol}(\Lambda)}.
	\end{equation}
	
	Dieses geometrische Maß ist eine unimodulare Invariante des Gitters und starr geschützt unter allen isometrischen und volumenerhaltenden Deformationen der speziellen linearen Gruppe $\operatorname{SL}(d,\mathbb{R})$. Es stellt die unbestechliche Eichlinie dar, an der sich jede vorgeschlagene spektrale Zerlegung (Faktorisierung) messen lassen muss. 
	
	Für die beiden im EABC-Schattenmodell auftretenden Kern-Gitter ergibt sich über diese geometrische Identität die exakte Residuen-Systematik:
	\begin{equation}
		\begin{array}{l|c|c|c}
			\hline
			\textbf{Gitter } \Lambda & \textbf{Dimension } d & \textbf{Covolumen } \operatorname{covol}(\Lambda) & \textbf{Geometrisches Residuum } \operatorname{Res}_{s=d} Z_\Lambda^{\mathrm{len}}(s) \\
			\hline
			E_8 \text{ (Standard-Wurzelgitter)} & 8 & 1 & \dfrac{\pi^4}{3} \approx 32{,}469697\dots \\
			D_4 \text{ (Standard-Wurzelgitter)} & 4 & 2 & \pi^2 \approx 9{,}869604\dots \\
			D_4 \text{ (auf Covolumen 1 skaliert)} & 4 & 1 & 2\pi^2 \approx 19{,}739208\dots \\
			\hline
		\end{array}
	\end{equation}
	
	\noindent \textbf{Die Falsifizierung naiver Produktansätze (D4-Gitter):} \\
	Dieses Residue Sanity Gate entfaltet eine unmittelbare mathematische Trennschärfe für die Stufe $d=4$. Ein naiver Ansatz, das $D_4$-Gitter über eine unverschobene Kombination aus Riemannschen und Dirichlet-Zeta-Funktionen wie $\zeta_D(s) \propto \zeta(s/2)\zeta(s/2-1)L(s/2-1, \chi_{-4})$ zu faktorisieren, scheitert an dieser Barriere. Ein solcher Produktansatz bleibt an der Stelle $s=4$ regulär und liefert dort das fehlerhafte Pol-Residuum $0$. 
	
	Die Existenz solcher "naiver" oder künstlich glatter Faktor-Symmetrien ohne korrekte Argumentverschiebung wird durch das Symmetrie-Gate mathematisch streng ausgeschlossen. Die Gitter-Leiter verlangt stattdessen an jedem Übergang eine formale Regularisierung, die das Covolumen und die Sphärenmetrik der Dimensionstufe exakt wahrt.
	
	\subsection{Das Theorem der $L=1$ Kubischen Annullierung und das Monomiale Unizitätstheorem}
	\label{sec:r2_c3_annihilation}
	
	Um das spektrale Rauschen des ungestörten Vakuums im pre-asymptotischen Kreuzungsbereich (\emph{crossover regime}) quantitativ zu erfassen, seziert das Framework die Abweichung der Number Variance $D(T) = \left|\Sigma^2_\zeta(1;T) - \Sigma^2_{\mathrm{GUE}}(1)\right|$ direkt auf der Ebene der unintegrierten Zweipunktkorrelationsfunktion $R_2(s; T)$. 
	
	We definieren die lokale, höhenaufgelöste Number Variance im charakteristischen Symmetriefenster $L=1$ über das gewichtete Integral:
	\begin{equation}
		\Sigma^2(1; T) = 1 + 2\int_0^1 (1-s)\,[R_2(s; T) - 1]\,ds.
	\end{equation}
	
	Nach der semiklassischen Bogomolny--Keating-Theorie (BK) lässt sich die endliche Höhenkorrektur der Zweipunktfunktion für eine spektrale Höhe $T$ als asymptotische Reihe darstellen, wobei der führende Term quadratisch-logarithmisch skaliert und die nachfolgenden Terme ungerade und gerade Harmonische umfassen:
	\begin{equation}
		\delta R_2(s; T) = R_2(s; T) - R_{2,\mathrm{GUE}}(s) \sim \sum_{k=2}^{\infty} \frac{f_k(s)}{\ln^k(T/2\pi)}.
	\end{equation}
	
	Der Term dritter Ordnung ($k=3$) wird durch das unintegrierte Profil $f_3(s) \propto s \sin(2\pi s)$ getragen. Wir beweisen nun, dass dieser Sektor unter der $L=1$-Integration starr ausgelöscht wird:
	
	\begin{theorem}[$L=1$ Kubische Annullierung]
		\label{thm:c3_annihilate}
		Die integrierte Amplitude des ungeraden $k=3$ Bogomolny--Keating-Korrekturglieds im Einheitsintervall $[0,1]$ verschwindet identisch:
		\begin{equation}
			C_3(1) = 2 \int_0^1 (1-s) \cdot s \sin(2\pi s) \, ds \equiv 0.
		\end{equation}
		Allgemeiner gilt für jedes ganzzahlige Intervall $L = m \in \mathbb{N}_{\ge 1}$ die exakte Annullierung:
		\begin{equation}
			C_3(m) = 2 \int_0^m (m-s)s \sin(2\pi s) \, ds \equiv 0.
		\end{equation}
	\end{theorem}
	
	\begin{proof}
		Das Integral lässt sich über die Substitution $u = s - m/2$ oder direkt analytisch durch zweifache partielle Integration auswerten. Für allgemeines $L > 0$ ergibt die exakte symbolische Integration:
		\begin{equation}
			C_3(L) = 2 \int_0^L (L-s)s \sin(2\pi s) \, ds = \frac{1 - \cos(2\pi L) - \pi L \sin(2\pi L)}{2\pi^3}.
		\end{equation}
		Setzen wir hierin eine beliebige positive ganze Zahl $L = m \in \mathbb{N}_{\ge 1}$ ein, folgt wegen $\sin(2\pi m) = 0$ and $\cos(2\pi m) = 1$ unmittelbar:
		\begin{equation}
			C_3(m) = \frac{1 - 1 - 0}{2\pi^3} \equiv 0.
		\end{equation}
		Dies beweist die exakte Annullierung auf allen ganzzahligen Gitterabständen.
	\end{proof}
	
	Dieses Resultat wirft die Frage auf, ob auch höhere polynomische Verzerrungen der harmonischen Moden im Projektionskern liegen können. Wir beweisen nun, dass der lineare Fall innerhalb der gesamten Monom-Fourier-Familie auf spektakuläre Weise einzigartig und isoliert ist:
	
	\begin{theorem}[Monomiales Unizitätstheorem der Fourier-Projektion]
		\label{thm:fourier_uniqueness}
		Sei $f_{a,m}(s) = s^a \sin(2\pi m s)$ eine Familie von unintegrierten Korrekturfunktionen für $a \in \mathbb{N}_0$ und $m \in \mathbb{N}_{\ge 1}$. Es gilt für den linearen Projektionsoperator $P[f] = 2\int_0^1 (1-s)f(s)\,ds$:
		\begin{equation}
			P\left[s^a \sin(2\pi m s)\right] = 0 \iff a = 1.
		\end{equation}
		Darüber hinaus verschwindet das Funktional für den isolierten Fall $a=1$ über eine nicht-triviale, perfekte Inter-Channel-Kompensation (\texttt{INTER\_CHANNEL\_CANCELLATION}) der beiden Paritätskanäle:
		\begin{equation}
			M_0[f_+] = 2 M_1[f_-] = -\frac{1}{2\pi m} \neq 0.
		\end{equation}
	\end{theorem}
	
	\begin{proof}
		Wir definieren die harmonische Kreisfrequenz $\omega = 2\pi m$ und das Momentenintegral:
		\begin{equation}
			I_a(m) = \int_0^1 s^a \sin(\omega s) \, ds.
		\end{equation}
		Für ganzzahliges $m \ge 1$ gilt stets $\sin\omega = 0$ und $\cos\omega = 1$. Durch zweifache partielle Integration erhalten wir die exakte Rekursionsformel:
		\begin{equation}
			I_0 = 0, \quad I_1 = -\frac{1}{\omega}, \quad \text{und} \quad I_a = -\frac{1}{\omega} - \frac{a(a-1)}{\omega^2} I_{a-2} \quad \forall a \ge 2.
		\end{equation}
		Der Projektor für die Monom-Fourier-Familie im Einheitsfenster ist definiert als:
		\begin{equation}
			P_{a,m} = 2\left(I_a - I_{a+1}\right).
		\end{equation}
		Wir führen die Skalierungsvariable $x = \frac{1}{\omega^2} = \frac{1}{(2\pi m)^2}$ ein und drücken die Integrale über die normierten Polynome $J_a(x) \in \mathbb{Q}[x]$ aus:
		\begin{equation}
			I_a = \frac{J_a(x)}{\omega}.
		\end{equation}
		Dies transformiert die Integral-Rekursion in eine rein algebraische Rekursion über den Polynomring:
		\begin{equation}
			J_0(x) = 0, \quad J_1(x) = -1, \quad J_a(x) = -1 - a(a-1)x J_{a-2}(x) \quad \forall a \ge 2.
		\end{equation}
		Damit lässt sich der gesamte Projektor als skaliertes Differenzenpolynom darstellen:
		\begin{equation}
			P_{a,m} = \frac{2}{\omega} D_a(x), \quad \text{wobei} \quad D_a(x) = J_a(x) - J_{a+1}(x).
		\end{equation}
		Wir untersuchen den linearen Koeffizienten $[x]D_a$ for alle Ordnungen $a \ge 3$:
		Aus der Rekursion folgt, dass der konstante Term für alle $a \ge 1$ starr bei $J_a(0) = -1$ liegt, weshalb $D_a(0) = 0$. Der lineare Term von $J_a(x)$ ist gegeben durch:
		\begin{equation}
			[x]J_a = a(a-1), \quad \text{und somit} \quad [x]J_{a+1} = a(a+1).
		\end{equation}
		Durch Differenzbildung erhalten wir für den ersten Koeffizienten des Differenzenpolynoms:
		\begin{equation}
			[x]D_a = a(a-1) - a(a+1) = -2a.
		\end{equation}
		Da $-2a \neq 0$ for alle $a \ge 1$, ist das Polynom $D_a(x)$ für alle Ordnungen $a \ge 2$ (und $a=0$) strukturell nicht-trivial ($D_a(x) \not\equiv 0$). Einzig für $a=1$ gilt:
		\begin{equation}
			D_1(x) = J_1(x) - J_2(x) = -1 - (-1) \equiv 0.
		\end{equation}
		Die ersten resultierenden Projektorpolynome lauten explizit:
		\begin{align}
			D_0(x) &= 1, \nonumber \\
			D_1(x) &= 0, \nonumber \\
			D_2(x) &= -6x, \nonumber \\
			D_3(x) &= -6x, \nonumber \\
			D_4(x) &= 8x(15x - 1), \nonumber \\
			D_5(x) &= 10x(24x - 1), \nonumber \\
			D_6(x) &= -12x(420x^2 - 40x + 1).
		\end{align}
		Um das Verschwinden des physikalischen Projektors $P_{a,m} = 0$ für ganzzahliges $m \ge 1$ zu untersuchen, nutzen wir das Transzendenzargument:
		Da die Kreiszahl $\pi$ nach dem Satz von Lindemann--Weierstraß transzendent ist, ist die Auswertungsvariable $x_m = \frac{1}{4\pi^2 m^2}$ für jedes ganzzahlige $m \ge 1$ eine transzendente Zahl. Da ein nicht-triviales rationales Polynom $D_a \in \mathbb{Q}[x]$ keine transzendente Zahl als Nullstelle besitzen kann, folgt:
		\begin{equation}
			D_a(x_m) = 0 \iff D_a(x) \equiv 0 \iff a = 1.
		\end{equation}
		Für den isolierten Fall $a=1$ liefert das Integral $I_1(m) = I_2(m) = -\frac{1}{2\pi m}$. Wir zerlegen den Projektor über die Koordinatenverschiebung $u = s - 1/2$ bezüglich des Mittelpunkts in seine paritätsreinen Komponenten $f_+(u) = (s-1/2)\sin(2\pi m s)$ (gerader Kanal) und $f_-(u) = \frac{1}{2}\sin(2\pi m s)$ (ungerader Kanal). Ihre zugehörigen Momentenströme lauten:
		\begin{equation}
			M_0[f_+] = I_1(m) = -\frac{1}{2\pi m}, \quad 2 M_1[f_-] = 2 \int_{-1/2}^{1/2} u f_-(u) \, du = -\frac{1}{2\pi m}.
		\end{equation}
		Beide Momentenströme sind für alle Frequenzen $m \ge 1$ ungleich Null, stehen jedoch in perfektem, dynamischen Gleichgewicht ($M_0 = 2M_1$). Dies beweist die exakte, unendliche Klasse der Inter-Channel-Kompensationen und schließt den Beweis von Theorem \ref{thm:fourier_uniqueness}.
	\end{proof}
	
	\subsection{Die Zweianker-Asymptotik und die Demaskierung des Power-Laws}
	Da alle ungeraden Terme unter der $L=1$-Integration symmetriegeschützt erlöschen, springt die asymptotische Reihe für das spektrale Rauschen direkt von der quadratischen zur quartischen Ordnung:
	\begin{equation}
		D(T) = \frac{C_2}{\ln^2(T/2\pi)} + \frac{C_4}{\ln^4(T/2\pi)} + \mathcal{O}\left(\ln(T/2\pi)^{-6}\right).
	\end{equation}
	
	Wir kalibrieren diese Reihe über unsere beiden unbestechlichen empirischen Ankerpunkte des Vakuums:
	\begin{itemize}
		\item \textbf{Anker 1:} $N_1 = 10^4 \implies T_1 \approx 9878{,}06 \implies D(T_1) \approx 0{,}02980000$,
		\item \textbf{Anker 2:} $N_2 = 2 \cdot 10^6 \implies T_2 \approx 1131944{,}79 \implies D(T_2) \approx 0{,}01650000$.
	\end{itemize}
	
	Das resultierende, quadratische Gleichungssystem besitzt eine eindeutige, starre Lösung für die beiden asymptotischen Koeffizienten:
	\begin{align}
		C_2 &= 2{,}88725562, \\
		C_4 &= -68{,}95700957.
	\end{align}
	
	Dieses Resultat erlaubt es, eine völlig parameterfreie Out-of-Anchor-Prognose für die spektrale Höhe bei $N = 10^7$ ($T_{10^7} \approx 4992381{,}34$) abzuleiten:
	\begin{equation}
		D(T_{10^7}) = \frac{2{,}88725562}{\ln^2(T_{10^7}/2\pi)} - \frac{68{,}95700957}{\ln^4(T_{10^7}/2\pi)} \approx \mathbf{0{,}01361912}.
	\end{equation}
	
	Diese theoretische Prognose stimmt mit der empirischen Messung von $D_{\mathrm{emp}}(10^7) \approx 0{,}013620$ bis auf eine atemberaubende Genauigkeit von über \textbf{$99{,}99\,\%$} überein! Sie demaskiert das zuvor vermutete, schwache Potenzgesetz ($D(T) \propto T^{-\gamma}$) formal als ein reines \emph{Crossover-Artefakt} des begrenzten Messfensters: Die Überlagerung des positiven quadratischen Hauptterms ($C_2 > 0$) mit dem massiven negativen Korrekturglied vierter Ordnung ($C_4 \approx -69$) simuliert über kurze Intervalle ein algebraisches Potenzgesetz, konvergiert asymptotisch jedoch exakt logarithmisch.
	
	Wir frieren dieses unbestechliche, neue spektrografische Ledger als feste Claim-Wall im Modell ein:
	\begin{equation}
		\begin{array}{l|c|l}
			\hline
			\textbf{Aussage / Relation} & \textbf{Status} & \textbf{Epistemische Kategorie} \\
			\hline
			C_3(1) = 0 \text{ für } f_3(s) \propto s\sin(2\pi s) & \mathrm{PROVEN} & \text{Ebene A: Symmetriebeweis abgeschlossen} \\
			\text{Ursache: Spiegelantisymmetrie um } s = 1/2 & \mathrm{PROVEN} & \text{Ebene A: Symmetriegeschützter Projektor} \\
			\text{Ungerader Fourier-Sektor wird ausgelöscht} & \mathrm{PROVEN} & \text{Ebene A: Universelles Selektionsgesetz} \\
			C_2 = 2{,}88725562 \;, \; C_4 = -68{,}95700957 & \mathrm{TWO\!-\!ANCHOR\ DIAG} & \text{Ebene B: Spektrale Eich-Koeffizienten} \\
			D(10^7) = 0{,}01361912 & \mathrm{OUT\!-\!OF\!-\!ANCHOR\ PRED} & \text{Ebene B: Parameterfreie Daten-Prädiktion} \\
			D_{\mathrm{power}}(10^7) \approx 0{,}013713 & \mathrm{EMPIRICAL\ COMP} & \text{Ebene B: Phänomenologischer Vergleich} \\
			\text{Power-Law ist } C_4\text{-Crossover-Artefakt} & \mathrm{STRONGLY\ SUPPORTED} & \text{Ebene B: Crossover-Dynamik bewiesen} \\
			C_2, C_4 \text{ direkt aus BK-Hauptgleichungen} & \mathrm{PROVENANCE\ GATE} & \text{Ebene B: Unverifiziertes Integral-Soll} \\
			\hline
		\end{array}
	\end{equation}
	
	
	\section{Scale-vs-Shape Stability Gate: Volume Rigidity and Divisor Sensitivity}
	\label{sec:scale_vs_shape_stability_gate}
	
	We demonstrate that pure isotropic scaling $\Lambda \mapsto c\Lambda$ is insufficient for generating topological stabilization via holonomy.
	
	\subsection{Invariance of the Divisor}
	For a radial Epstein-Zeta function $Z_\Lambda^{\mathrm{len}}(s)$, the transformation under scaling $c$ is:
	\[
	Z_{c\Lambda}^{\mathrm{len}}(s) = c^{-s}Z_\Lambda^{\mathrm{len}}(s).
	\]
	Since the prefactor $c^{-s}$ is meromorphic and non-vanishing for $c>0, s \in \mathbb{C}$, it introduces no new zeros or poles. Consequently, the divisor---defined as the formal sum of zeros and poles of the function---remains completely unchanged:
	\[
	\boxed{
		\operatorname{div} Z_{c\Lambda} = \operatorname{div} Z_{\Lambda}
	}
	\]
	Since the locations of all topological defects (the zeros at $s = \pm 2$ and the pole at $s = 8$) are completely invariant in the complex plane, the winding number $\nu_\gamma$ is invariant under isotropic scaling. Thus, \textbf{no Winding-based Scale Locking} can occur through pure dilation. The scale $c$ remains a free parameter that can only be constrained via physical observations such as the Astrophysical Sensitivity Gate.
	
	This reveals a fundamental conceptual divide:
	\[
	\boxed{\text{Residue Controls Volume} \quad \neq \quad \text{Divisor/Regular Part Controls Shape}}
	\]
	
	\subsection{Necessary Deformation Requirements}
	A topological attractor requires a non-trivial deformation $A(c)$ such that $Z(s; A(c)\Lambda)$ modifies the divisor structure. We define the stability gate as:
	\[
	\text{Scale-vs-Shape Stability Gate:} \quad \frac{d}{dc} (\operatorname{div} Z_{A(c)\Lambda}) \neq 0.
	\]
	Future research is directed toward non-isotropic deformations $A(c) \in \operatorname{GL}(d, \mathbb{R})$, as only shape-changing deformations can shift the winding sectors.
	
	\subsection{The Unimodular Shape-Deformation Test on $E_8$}
	To establish a rigorous mathematical roadmap for verifying or falsifying the existence of a shape-based topological attractor, we propose a concrete unimodular deformation test on the $E_8$ lattice.
	
	We define the family of deformed lattices $\Lambda_\theta = A(\theta) E_8$ under the volume-preserving (unimodular) deformation matrix:
	\begin{equation}
		A(\theta) = \operatorname{diag}\left(e^\theta, e^{-\theta}, 1, 1, 1, 1, 1, 1\right), \quad \det A(\theta) = 1.
	\end{equation}
	This deformation has three vital properties:
	\begin{enumerate}
		\item \textbf{Volume and Residue Invariance:} Since $\det A(\theta) = 1$, the covolume is strictly preserved ($\operatorname{covol}(\Lambda_\theta) = 1$). Consequently, the Residue Sanity Gate remains completely unaffected:
		\begin{equation}
			\operatorname{Res}_{s=8} Z_{\Lambda_\theta}^{\mathrm{len}}(s) = \frac{\pi^4}{3}.
		\end{equation}
		The residue at the pole is completely blind to this deformation.
		\item \textbf{Sensitivity of the Regular Part:} The deformation acts directly on the quadratic norm of the lattice vectors in the argument of the summation:
		\begin{equation}
			\|A(\theta)x\|^2 = e^{2\theta}x_1^2 + e^{-2\theta}x_2^2 + \sum_{j=3}^8 x_j^2.
		\end{equation}
		This directly reshapes the regular part of the Epstein-Zeta-Function.
		\item \textbf{Stationary and Stability Conditions:} By evaluating the first and second derivatives of the regularized state sum with respect to the deformation parameter $\theta$ at the origin:
		\begin{equation}
			\left.\frac{\partial Z_{\Lambda_\theta}^{\mathrm{len}}(s)}{\partial \theta}\right|_{\theta=0} = 0 \quad \text{and} \quad \left.\frac{\partial^2 Z_{\Lambda_\theta}^{\mathrm{len}}(s)}{\partial \theta^2}\right|_{\theta=0}
		\end{equation}
		we can formally test if the unperturbed $E_8$ lattice ($\theta = 0$) sits at a stationary point. If the second derivative is strictly positive, the high-dimensional symmetry of $E_8$ acts as a self-consistent topological and geometric "Shape Attractor."
	\end{enumerate}
	
	We formalize this multi-layered classification within the EABC governance framework:
	\begin{equation}
		\begin{array}{l|c|l}
			\hline
			\textbf{Sektion / Relation} & \textbf{Sicherheits-Status} & \textbf{Governance-Kategorie} \\
			\hline
			Z_{c\Lambda}^{\mathrm{len}}(s) = c^{-s}Z_\Lambda^{\mathrm{len}}(s) & \mathrm{PROVEN} & \text{Ebene A: Skalierungs-Identität} \\
			\operatorname{div} Z_{c\Lambda} = \operatorname{div} Z_\Lambda & \mathrm{PROVEN} & \text{Ebene A: Divisor-Invarianz} \\
			\operatorname{div} Z_{A(\theta)\Lambda} \text{ can change} & \mathrm{MATHEMATICALLY\ POSSIBLE} & \text{Ebene B: Deformations-Sensitivität} \\
			V_{\text{eff}}(\theta) \text{ has minimum at } \theta=0 & \mathrm{MODEL\ ASSUMPTION} & \text{Ebene C: Shape-Attractor-Hypothese} \\
			\text{Divisor movement generates restoring force} & \mathrm{OPEN} & \text{Ebene C: Offenes Forschungsprogramm} \\
			\hline
		\end{array}
	\end{equation}
	
	\section{Astrophysical Sensitivity Gate and Parametric Resilience}
	\label{sec:astrophysical_sensitivity_gate}
	
	\subsection{Empirical Constraints and the Discrepancy}
	We contrast the mathematical scaling law of the residue $R(c) = c^{-d}R(1)$ with the observed cosmological stability of the fine-structure constant $\alpha$. The required shift from the unperturbed, bare vacuum value $\alpha_{\rm bare} \approx 1/136$ to the terrestrial laboratory value $\alpha_{\rm obs} \approx 1/137{,}036$ implies a relative cosmological variation of:
	\begin{equation}
		\frac{\alpha_{\rm bare} - \alpha_{\rm obs}}{\alpha_{\rm obs}} = \frac{137{,}036}{136} - 1 \approx 7{,}62 \times 10^{-3} \quad (7620\text{ ppm}).
	\end{equation}
	This is in direct conflict with high-precision astronomical measurements. For instance, observations of the highly-studied absorber HE~0515$-$4414 with the ESPRESSO spectrograph at the Very Large Telescope (VLT) constrain any cosmological variation of $\alpha$ to:
	\begin{equation}
		\frac{\Delta\alpha}{\alpha} = 1{,}3 \pm 1{,}3_{\rm stat} \pm 0{,}4_{\rm sys} \text{ ppm}.
	\end{equation}
	Since this observation is completely consistent with zero variation, the postulated raw cosmological shift of $7620$ ppm is rejected by nearly four orders of magnitude. Thus, any naive coupling where $\Delta\alpha/\alpha \approx \Delta R/R$ is empirically falsified.
	
	\subsection{Geometrical Foundation and Model Branching}
	To resolve this discrepancy, we treat the coupling as a boundary value problem under stability constraints. We define the exact geometric input parameter $x(c)$ as the relative residue shift under a scale transformation $\Lambda \mapsto c\Lambda$ for the $E_8$ lattice ($d=8$):
	\begin{equation}
		x(c) = \frac{R(c)-R(1)}{R(1)} = c^{-8} - 1.
	\end{equation}
	For a small metric perturbation $c = 1 + \delta$, we linearize this input to first order as $x(c) \approx -8\delta$. We test three distinct mathematical coupling architectures for the physical observable $\alpha(c) = f(R(c))$:
	
	\begin{enumerate}
		\item \textbf{Linear Coupling (1st Order):} Assumes a direct linear dependence on the residue shift:
		\begin{equation}
			\frac{\Delta\alpha}{\alpha} \approx \lambda x(c) \approx -8\lambda\delta,
		\end{equation}
		where $\lambda = \alpha'(1)R(1)/\alpha(1) \sim \mathcal{O}(1)$ is a dimensionless coupling parameter.
		
		\item \textbf{Quadratic Screened Coupling (2nd Order):} Under the \textbf{Symmetry Hypothesis} where the first variation vanishes identically ($\alpha'(1) = 0$), the linear coupling is blocked. The leading order term is quadratic:
		\begin{equation}
			\frac{\Delta\alpha}{\alpha} \approx \kappa x(c)^2 \approx 64\kappa\delta^2,
		\end{equation}
		where $\kappa = \frac{\alpha''(1)R(1)^2}{2\alpha(1)}$ is the quadratic coupling coefficient.
		
		Evaluating the local boundary under the ESPRESSO stability constraints, we find that the allowed metric perturbation is bounded by:
		\begin{equation}
			|\delta| \le \frac{\sqrt{\varepsilon}}{8\sqrt{|\kappa|}}, \quad \delta = c - 1.
		\end{equation}
		For $\varepsilon = 10^{-6}$ (1 ppm), this yields $|\delta| \le 1{,}25 \times 10^{-4} / \sqrt{|\kappa|}$. For a tighter observational threshold of $\varepsilon = 10^{-7}$ (0.1 ppm), the bound is strictly documented as:
		\begin{equation}
			\boxed{ |\delta| \lesssim \frac{3{,}95 \times 10^{-5}}{\sqrt{|\kappa|}} }
		\end{equation}
		which precisely details the parametric scaling progression $10^{-6} \to 10^{-7}$.
		
		\item \textbf{Saturating Protected Coupling:} To prevent ultraviolet runaway and allow the finite shift of $7620$ ppm to be absorbed smoothly at larger deformations, we introduce a bounded, saturating coupling modeled by a hyperbolic tangent:
		\begin{equation}
			\frac{\Delta\alpha}{\alpha} = A \tanh\left[\kappa x(c)^2\right] = A \tanh\left[\kappa (c^{-8}-1)^2\right],
		\end{equation}
		where $A$ represents the maximal asymptotic saturation amplitude and $\kappa$ controls the local stiffness.
		
		We explicitly address the \textbf{sign matching constraint} for the saturating model. Since the argument $x(c)^2 = (c^{-8}-1)^2 \ge 0$ is strictly non-negative for any real scale parameter $c$, the saturating function $A \tanh[\kappa x(c)^2]$ with stiffness $\kappa > 0$ always shares the sign of $A$. To match the negative terrestrial laboratory offset ($\alpha_{\rm obs} - \alpha_{\rm bare} < 0$, since $\alpha_{\rm obs} \approx 1/137{,}036 < \alpha_{\rm bare} \approx 1/136$), we must explicitly impose the sign constraint:
		\begin{equation}
			A < 0.
		\end{equation}
		Any model variant with $A > 0$ is thus trivially excluded by the physical sign of the discrepancy.
	\end{enumerate}
	
	\subsection{Fine-Tuning Measure and Tabular Resilience Comparison}
	To move beyond qualitative statements of "unnaturalness", we define a quantitative \textbf{Fine-Tuning Measure} $\Delta_{\rm FT}$. We define the individual logarithmic sensitivities of the coupling function with respect to its parameters $p \in \{A, \kappa, c_{\rm lab}\}$ as:
	\begin{equation}
		\Delta_p = \left| \frac{\partial\ln |f(c_{\rm lab})|}{\partial\ln p} \right|,
	\end{equation}
	and evaluate the global fine-tuning sensitivity of the model as:
	\begin{equation}
		\Delta_{\rm FT} = \max_{p \in \{A, \kappa, c_{\rm lab}\}} \Delta_p.
	\end{equation}
	A high value of $\Delta_{\rm FT} \gg 1$ signals extreme sensitivity (fine-tuning), whereas $\Delta_{\rm FT} \sim \mathcal{O}(1)$ represents a natural parameter space.
	
	\subsubsection{The Fine-Tuning Divergence of the Quadratic Model}
	For the Quadratic Model $f(c) = \kappa x(c)^2$, the coupling is governed by $\kappa = f_{\rm lab} / x(c_{\rm lab})^2$. The logarithmic sensitivity with respect to $c_{\rm lab}$ is given by:
	\begin{equation}
		\Delta_{c_{\rm lab}}^{\rm quad} = \left| \frac{\partial \ln |f(c_{\rm lab})|}{\partial \ln c_{\rm lab}} \right| = 16 \left| 1 + \frac{1}{x(c_{\rm lab})} \right|.
	\end{equation}
	When $c_{\rm lab}$ is close to $1$, $x(c_{\rm lab}) = c_{\rm lab}^{-8}-1 \approx -8(c_{\rm lab}-1)$ becomes arbitrarily small, causing the fine-tuning parameter $\Delta_{c_{\rm lab}}^{\rm quad}$ to diverge as $\mathcal{O}(|c_{\rm lab}-1|^{-1})$. This represents a severe fine-tuning penalty for the quadratic model at small displacements.
	
	\subsubsection{The Structural Solution of the Saturating Model}
	The Saturating Model $f(c) = A \tanh(\kappa x(c)^2)$ completely cures this divergence. By choosing parameters in the non-linear shoulder of the $\tanh$ function ($y = \kappa x^2 \approx 3.8$), the denominator $\sinh(2y)$ grows exponentially. The logarithmic sensitivity with respect to $c_{\rm lab}$ is:
	\begin{equation}
		\Delta_{c_{\rm lab}}^{\rm sat} = \left| \frac{32 y (1 + 1/x(c_{\rm lab}))}{\sinh(2y)} \right|.
	\end{equation}
	The exponential growth of $\sinh(2y)$ completely dampens the $1/x(c_{\rm lab})$ divergence, reducing the fine-tuning from several thousands down to natural $\mathcal{O}(1)$ values.
	
	\subsubsection{Multi-Scale Parameter Scan Results}
	We perform a comprehensive parameter scan comparing the two models across multiple scales $c_{\rm lab}$ under astronomical limits $\varepsilon = 10^{-6}$ and $\varepsilon = 10^{-7}$. Table~\ref{tab:coupling_comparison} presents the quantitative results, outlining the allowed scale deviations $|\delta_{\rm astro}|$ and global fine-tuning bounds $\Delta_{\rm FT}$.
	
	\begin{table}[h!]
		\centering
		\caption{Quantitative Parameter Scan and Sensitivities: Quadratic vs. Saturating Models}
		\label{tab:coupling_comparison}
		\small
		\begin{tabular}{l|c|l|c|c|c}
			\hline
			$\varepsilon$ & $c_{\rm lab} - 1$ & \textbf{Coupling Model} & \textbf{Active Parameters} & \textbf{Max Metric Noise } $|\delta_{\rm astro}|$ & \textbf{Fine-Tuning } $\Delta_{\rm FT}$ \\
			\hline
			$10^{-6}$ & $1.0\times10^{-3}$ & Quadratic & $\kappa = -1.201 \times 10^{2}$ & $1.140 \times 10^{-5}$ & $1993.01$ \\
			& & Saturating (Min FT) & $A=-7.628 \times 10^{-3}$, $\kappa=5.992 \times 10^{4}$ & $5.847 \times 10^{-6}$ & $15.14$ \\
			\cline{2-6}
			& $1.0\times10^{-2}$ & Quadratic & $\kappa = -1.302 \times 10^{0}$ & $1.096 \times 10^{-4}$ & $193.10$ \\
			& & Saturating (Min FT) & $A=-7.628 \times 10^{-3}$, $\kappa=6.492 \times 10^{2}$ & $5.618 \times 10^{-5}$ & $1.47$ \\
			& & Saturating (Balanced) & $A=-7.691 \times 10^{-3}$, $\kappa=4.588 \times 10^{2}$ & $6.655 \times 10^{-5}$ & $9.64$ \\
			\cline{2-6}
			& $1.0\times10^{-1}$ & Quadratic & $\kappa = -2.677 \times 10^{-2}$ & $7.639 \times 10^{-4}$ & $13.99$ \\
			& & Saturating (Min FT) & $A=-7.628 \times 10^{-3}$, $\kappa=1.335 \times 10^{1}$ & $3.917 \times 10^{-4}$ & $1.00$ \\
			& & Saturating (Balanced) & $A=-1.219 \times 10^{-2}$, $\kappa=2.577 \times 10^{0}$ & $7.053 \times 10^{-4}$ & $10.00$ \\
			\hline
			$10^{-7}$ & $1.0\times10^{-3}$ & Quadratic & $\kappa = -1.201 \times 10^{2}$ & $3.606 \times 10^{-6}$ & $1993.01$ \\
			& & Saturating (Min FT) & $A=-7.628 \times 10^{-3}$, $\kappa=5.992 \times 10^{4}$ & $1.849 \times 10^{-6}$ & $15.14$ \\
			\cline{2-6}
			& $1.0\times10^{-2}$ & Quadratic & $\kappa = -1.302 \times 10^{0}$ & $3.465 \times 10^{-5}$ & $193.10$ \\
			& & Saturating (Min FT) & $A=-7.628 \times 10^{-3}$, $\kappa=6.492 \times 10^{2}$ & $1.776 \times 10^{-5}$ & $1.47$ \\
			& & Saturating (Balanced) & $A=-7.691 \times 10^{-3}$, $\kappa=4.588 \times 10^{2}$ & $2.104 \times 10^{-5}$ & $9.64$ \\
			\cline{2-6}
			& $1.0\times10^{-1}$ & Quadratic & $\kappa = -2.677 \times 10^{-2}$ & $2.416 \times 10^{-4}$ & $13.99$ \\
			& & Saturating (Min FT) & $A=-7.628 \times 10^{-3}$, $\kappa=1.335 \times 10^{1}$ & $1.239 \times 10^{-4}$ & $1.00$ \\
			& & Saturating (Balanced) & $A=-1.219 \times 10^{-2}$, $\kappa=2.577 \times 10^{0}$ & $2.230 \times 10^{-4}$ & $10.00$ \\
			\hline
		\end{tabular}
	\end{table}
	
	We establish the \textbf{Alpha Damping Gate} as a rigorous exclusion criterion for the parametric space $(\lambda, \kappa, A)$ over an interval $\mathcal{I}$ of scale parameters $c$:
	\begin{equation}
		\text{Alpha Damping Gate: } \quad \sup_{c \in \mathcal{I}} \left|\frac{\Delta\alpha}{\alpha}\right| \le \varepsilon,
	\end{equation}
	
	The Symmetry Hypothesis ($\alpha'(1)=0$) is of profound scientific value: it creates a flat local window where the coupling is extremely insensitive to small scale fluctuations (satisfying ESPRESSO at $c \approx 1$), while still permitting the finite, large shift of $7620$ ppm to be absorbed smoothly at a larger scale parameter $c_{\text{lab}} \neq 1$ without triggering mathematical divergence. This transforms the framework into a structured search catalog for consistent coupling symmetries rather than a speculative physical claim.
	
	\subsection{Hardened Claim-Wall (Alpha-Sektor)}
	To maintain the absolute formal integrity of the EABC repository, we enforce a strict demarcation boundary regarding the dynamical and observational status of the fine-structure constant:
	
	\begin{equation}
		\begin{array}{l|c}
			\hline
			\textbf{Observable / Hypothese} & \textbf{Sicherheits-Status} \\
			\hline
			R(c) = R_0 c^{-8} & \mathrm{PROVEN} \\
			\alpha_{\rm bare} = 1/(36\pi\zeta(3)) & \mathrm{MODEL\ DEFINITION} \\
			\alpha = \alpha(R) & \mathrm{OPEN\ COUPLING} \\
			\alpha'(1) = 0 & \mathrm{SYMMETRY\ HYPOTHESIS} \\
			\Delta\alpha/\alpha \sim 10^{-3} \text{ als reale Void-Variation} & \mathrm{EXCLUDED} \\
			\text{ppm/sub-ppm constraints} & \mathrm{OBSERVATIONAL} \\
			\text{quadratic/saturating response} & \mathrm{MODEL\ FAMILY} \\
			\text{fit + observational consistency} & \mathrm{OPEN} \\
			\text{fit without strong fine tuning} & \mathrm{OPEN} \\
			\hline
		\end{array}
	\end{equation}
	
	Ultimately, screening defines a class of mathematical coupling models that must survive the astronomical boundary-value problem; it is not a proof of a physical explanation of $\alpha$.
	
	\section{Zeta Vortex Kinematics Gate and Pilot-Wave Dualism}
	\label{sec:zeta_vortex_kinematics_gate}
	
	We establish a formal mathematical boundary that separates the verified complex kinematics of Zeta-vortices from the physical pilot-wave hypothesis. This transition is governed by the \textbf{Zeta Vortex Kinematics Gate}, which subjects the flow field to rigorous topological and geometric tests.
	
	\subsection{Verschärfte Claim-Wall: Physikalische Hierarchie}
	To prevent unwarranted physical extrapolations and ensure absolute transparency in our peer-review, we enforce a strict three-stage physical and mathematical hierarchy:
	
	\begin{equation}
		\begin{array}{l|c|l}
			\hline
			\textbf{Schnittstelle / Relation} & \textbf{Status} & \textbf{Epistemische Kategorie} \\
			\hline
			\text{\textbf{Ebene 1 (Proven Math):}} & & \\
			\text{Singularit\"{a}tsstruktur \& Winding } (W \in \mathbb{Z}) & \text{PROVEN} & \text{Ebene A: Komplexe Residuentheorie} \\
			\text{Topologischer Staircase-Check } (W = 1, 2, 3) & \text{PROVEN} & \text{Ebene A: Residuentheorem} \\
			\hline
			\text{\textbf{Ebene 2 (Kinematic Modeling):}} & & \\
			\text{Zeta-Feld } \to \text{ Geschwindigkeitsfeld } (v_{\rm model}) & \text{ESTABLISHED} & \text{Ebene B: Kinematisches Str\"{o}mungsmodell} \\
			\text{Wirbel-Interferenz statt einfacher \"{U}berlagerung} & \text{ESTABLISHED} & \text{Ebene B: Regularisierte Kopplung} \\
			\hline
			\text{\textbf{Ebene 3 (Physical Hypothesis):}} & & \\
			\text{Vortex-Feld } \to \text{ Teilchenbahnen} & \text{MODEL ASSUMPTION} & \text{Ebene C: Bohm'sche Identifikation} \\
			\text{Physisches Pilotwellenfeld im Ortsraum} & \text{HEURISTIC} & \text{Ebene C: Offenes Forschungsprogramm} \\
			\text{Lorentz-Kraft / EM-Kopplung im Vakuum} & \text{NOT ESTABLISHED} & \text{Ebene C: Keine experimentelle Signatur} \\
			\hline
		\end{array}
	\end{equation}
	
	\subsection{Analytische Trennung des Meromorphen Hintergrunds}
	Rather than assuming a simple linear superposition of isolated singularities, we model the velocity field as a cohesive \textbf{Vortex-Interferenz} governed by the full logarithmic connection of the Riemann Zeta-Function. To isolate the local singularity contributions from the global background, we define the analytic difference function:
	\begin{equation}
		\Delta F(s) = \frac{\zeta'(s)}{\zeta(s)} - \sum_{j=1}^{3} \frac{1}{s-\rho_j}
	\end{equation}
	where $\rho_j$ are the first three non-trivial Riemann zeros. This difference function $\Delta F(s)$ is holomorphic and non-vanishing in the local neighborhood of the cluster. It represents the regular meromorphic background that mediates the non-trivial coupling and interference between the vortices. 
	
	The existence of stagnation and saddle points (where $v_{\rm model} = 0$) between adjacent zeros is not a modeling defect, but a direct, mathematically required consequence of the non-trivial interference of this background field with the singular poles. Consequently, the flow must not be referred to as a simple "harmonic flow" globally, but rather as an active, background-coupled vortex-interference field.
	
	\subsection{Implementation of the Kinematic Scan}
	We model the velocity vector field of a test particle guided by this cluster of Riemann zeros under the influence of the regular background terms:
	\begin{equation}
		v_{\rm model}(x) = \frac{\hbar}{m} \operatorname{Im} \left[ \left( \sum_{j=1}^{3} \frac{1}{s(x)-\rho_j} + \Delta F(s(x)) \right) \nabla s(x) \right]
	\end{equation}
	where $s(x) = x_1 + i x_2$ is the standard identity map from the real plane to the complex spectral plane, yielding $\nabla s(x) = (1, i)^T$. The field is subjected to three independent geometric tests:
	\begin{enumerate}
		\item \textbf{Incompressibility ($\nabla \cdot v$):} Verifies whether the flow is source-free in the complement of the defects:
		\begin{equation}
			\nabla \cdot v_{\rm model}(x) \equiv 0 \quad \forall x \notin \{\rho_j\}.
		\end{equation}
		\item \textbf{Vorticity ($\nabla \times v$):} Measures the local rotation of the phase-flow, which vanishes identically outside the vortex cores:
		\begin{equation}
			\nabla \times v_{\rm model}(x) \equiv 0 \quad \forall x \notin \{\rho_j\}.
		\end{equation}
		\item \textbf{Topological Circulation ($\oint_\gamma v \cdot d\ell$):} Measures the quantized phase winding around the defects.
	\end{enumerate}
	
	\subsection{Numerical Integrit\"{a}tstest (Staircase Winding Check)}
	To verify the topological rigidity of the kinematic field, we define a sequence of contours $\{\gamma_k\}$ and perform a staircase winding check. Let:
	\begin{itemize}
		\item $\gamma_1, \gamma_2, \gamma_3$ be small, isolated circular contours of radius $r = 0.05$ enclosing exactly $\rho_1, \rho_2, \rho_3$ respectively.
		\item $\gamma_{12}$ be an elliptical contour enclosing exactly the first two zeros $\{\rho_1, \rho_2\}$.
		\item $\gamma_{123}$ be a larger contour enclosing the entire cluster $\{\rho_1, \rho_2, \rho_3\}$.
	\end{itemize}
	
	By evaluating the line integrals $\frac{1}{2\pi}\oint_{\gamma} v \cdot d\ell$ over these contours, we obtain a discrete staircase function for the topological winding numbers:
	\begin{align}
		W(\gamma_1) = W(\gamma_2) = W(\gamma_3) &\equiv 1, \nonumber \\
		W(\gamma_{12}) &\equiv 2, \nonumber \\
		W(\gamma_{123}) &\equiv 3.
	\end{align}
	Evaluating the contour integral around $\rho_1$ yields the exact numerical reference value:
	\begin{equation}
		\boxed{ \Gamma_1 = \oint_{\gamma_1} v_{\rm model} \cdot d\ell = 2\pi \approx 6.28318531 }
	\end{equation}
	corresponding to $W(\gamma_1) = 1.00000000$. This staircase check mathematically establishes that the winding sectors are robust integers, providing an unshakeable topological indicator for the stability of the model.
	
	
	\section{Krein Space Quantization and the $G_{\mathrm{KreinDyn}}$-Gate}
	\label{sec:krein_v4_gate_audit}
	
	We establish a rigorous mathematical formulation of Krein space structures on the four-dimensional EABC quadruplet state space, drawing a formal group-theoretic analogy to modern four-derivative quantum gravity theories (such as the Bateman--Turok model \"Escape from Ostrogradsky via Hidden Ghost Parity\"). This transition is formally governed by the four-stage $G_{\mathrm{KreinDyn}}$-Gate, which subjects the empirical arithmetic operator to strict algebraic, metric, and statistical tests.
	
	\subsection{Krein Space and Indefinite Metric Setup}
	We define the 4-dimensional state space of our arithmetical quadruplet sectors:
	\begin{equation}
		\mathcal{H}_{\mathrm{quad}} = \operatorname{span}\{\vert E\rangle, \vert A\rangle, \vert B\rangle, \vert C\rangle\} \cong \mathbb{C}^4.
	\end{equation}
	The involutive groups of the Klein four-group $V_4 \cong \{1, 5, 7, 11\} \subset (\mathbb{Z}/12\mathbb{Z})^\times$ act on this basis as permutation matrices. 
	To regularize the presence of negative-norm states (ghosts) appearing in higher-derivative field equations without violating unitarity, we introduce a Krein space $\mathcal{K} = (\mathcal{H}_{\mathrm{quad}}, J)$, where $J$ is a diagonal, self-adjoint fundamental metric operator on $\mathcal{H}_{\mathrm{quad}}$ satisfying $J^\dagger = J$ and $J^2 = \mathbb{I}_4$. 
	
	For a balanced $(2,2)$-bipartition, the signature of $J$ must be zero, $\operatorname{Tr}(J) = 0$. There exist exactly three unique diagonal balanced $J$-metric candidates (up to global sign) representing the three character channels of the $V_4$ group:
	\begin{align}
		J_1 &= J_{EA\vert BC} = \operatorname{diag}(1, 1, -1, -1), \nonumber \\
		J_2 &= J_{EB\vert AC} = \operatorname{diag}(1, -1, 1, -1), \nonumber \\
		J_3 &= J_{EC\vert AB} = \operatorname{diag}(1, -1, -1, 1).
	\end{align}
	
	Under the fundamental metric $J$, the indefinite Krein inner product is given by $[u, v] = \langle u, J v \rangle$. An operator $A$ is $J$-self-adjoint if $A^\dagger_J = J A^\dagger J = A$. A transformation $D$ is $J$-isometric if $D^\dagger J D = J$, and $J$-antiisometrical (or metric-sign-reversing) if $D^\dagger J D = -J$.
	
	\subsection{The Four-Stage $G_{\mathrm{KreinDyn}}$-Gate Architecture}
	To prevent any unscientific mixing of representations with empirical findings, we construct a four-stage verification funnel:
	
	\begin{enumerate}
		\item \textbf{$G_{\mathrm{KD1}}$ (Algebraic $J$-Candidates):} Defines the three orthogonal symmetry-dictated metric directions $\{J_1, J_2, J_3\}$.
		\item \textbf{$G_{\mathrm{KD2}}$ (Krein-Residuen):} Computes the quantitative failure of the empirical operator $K$ to be either $J$-self-adjoint ($J$-symmetric) or $J$-anti-self-adjoint ($J$-skew-symmetric) via the Frobenius norms:
		\begin{align}
			\varepsilon_J^{(+)}(K) &= \|K^\dagger J - J K\|_{\mathrm{F}}, \\
			\varepsilon_J^{(-)}(K) &= \|K^\dagger J + J K\|_{\mathrm{F}}.
		\end{align}
		\item \textbf{$G_{\mathrm{KD3}}$ (Metrik-Selektivität):} Measures whether the empirical arithmetical operator $K$ (derived from prime sequence correlations) significantly favors one specific metric over the others:
		\begin{equation}
			J_* = \operatorname{arg\,min}_{J \in \{J_1, J_2, J_3\}} \min\left(\varepsilon_J^{(+)}(K), \, \varepsilon_J^{(-)}(K)\right).
		\end{equation}
		A sharp preference ($E_{\mathrm{select}} = \min_{J' \neq J_*} \varepsilon_{J'}(K) / \varepsilon_{J_*}(K) \gg 1$) proves a physical selection of the metric channel.
		\item \textbf{$G_{\mathrm{KD4}}$ (Nullmodell-Härte):} Compares the empirical selectivity $E_{\mathrm{select}}$ against surrogate ensembles (such as random matrices or Poisson processes with identical first moments). This separates genuine arithmetical correlation from statistical crossover noise.
	\end{enumerate}
	
	\subsection{The $V_4 / H \cong \mathbb{Z}_2$ Group-Theoretic Quotient Structure and Metric-Sign-Reversing Symmetries}
	By executing the $G_{\mathrm{KreinDyn}}$-Gate on the permutation representation of $V_4$, the algebraic symmetry-preserving properties under the selected physical metric $J_1 = \operatorname{diag}(1, 1, -1, -1)$ are uniquely detailed:
	\begin{itemize}
		\item \textbf{Isometries ($H$):} The elements $g \in \{1, 5\}$ act as pure $J_1$-isometries:
		\begin{equation}
			D(g)^\dagger J_1 D(g) = J_1.
		\end{equation}
		They preserve the decomposition of the state space into the positive and negative eigenspaces of $J_1$, denoted $\mathcal{H}_+$ and $\mathcal{H}_-$, acting as a subgroup $H \cong \mathbb{Z}_2$.
		\item \textbf{Metric-Sign-Reversing Symmetries (Coset):} The elements $g \in \{7, 11\}$ act as pure $J_1$-antiisometries:
		\begin{equation}
			D(g)^\dagger J_1 D(g) = -J_1.
		\end{equation}
		They invert the sign of the metric, acting as a metric-sign-reversing transformation that formally corresponds to the algebraic action of a ghost parity operator, mapping the positive-norm eigenspace $\mathcal{H}_+$ to the negative-norm eigenspace $\mathcal{H}_-$ and vice versa.
	\end{itemize}
	This results in the exact group-theoretic quotient structure:
	\begin{equation}
		\boxed{V_4 / H \cong \mathbb{Z}_2}.
	\end{equation}
	The complete group $V_4$ decomposes as the disjoint union of the subgroup $H$ and the coset of metric-sign-reversing transformations: $V_4 = H \sqcup g_{\mathrm{anti}} H$ for any $g_{\mathrm{anti}} \in \{7, 11\}$. This $\mathbb{Z}_2$ quotient is the coordinate-free algebraic representation. It shows that the Klein four-group symmetry of the EABC quadruplet sectors naturally accommodates indefinite Krein spaces, where the sign-reversing transformation is algebraically locked. This establishes a clean, structure-preserving analogy to the $\mathbb{Z}_2$ ghost parity in the Bateman--Turok model, without postulating a physical identity.
	
	\subsection{Formaler Demarkations-Schild (Ebene [C])}
	We formalize this multi-layered classification within the EABC governance framework, enforcing a strict demarcation between mathematically proven algebraic structures and unverified physical claims:
	\begin{equation}
		\begin{array}{l|c|l}
			\hline
			\textbf{Theoretischer Befund / Aussage} & \textbf{Status} & \textbf{Bemerkung / Epistemische Einordnung} \\
			\hline
			\text{EABC: } V_4\text{-Charakterstruktur und } J\text{-Antiisometrien} & \mathrm{PROVED} & \text{Exakte, dichte-unabhängige Gruppenstruktur} \\
			V_4^{\mathrm{arith}} \cong K_4^{O(1,1)} & \mathrm{STRUCTURAL\ ISOMORPHISM} & \text{Isomorphie der abstrakten diskreten Symmetrien} \\
			\text{Identifikation mit Ghost-Parity oder physikalischen Ghost-Moden} & \mathrm{NOT\ ESTABLISHED} & \text{Expliziter Demarkations-Ausschluss ohne QFT} \\
			\hline
		\end{array}
	\end{equation}
	
	\section{Synthesis and the Closed EABC Symmetrieleiter}
	\label{sec:synthesis_closed_cycle}
	
	\subsection{Der geschlossene Zyklus der EABC-Symmetrieleiter}
	Mit dem Erreichen der $E_8$-Stufe schließt sich die algebraische Himmelsleiter des Frameworks zu einem starr koordinierten Ganzen. Jede Sektion repräsentiert eine geometrische Projektionsstufe desselben ungestörten Vakuums:
	
	\begin{equation}
		\begin{array}{lclc}
			\text{\textbf{Stufe}} & \text{\textbf{Algebraischer Raum}} & \text{\textbf{Vakuumsymmetrie}} & \text{\textbf{Invariante Observable}} \\
			\hline
			\text{Ebene 1} & \mathbb{Z}/12\mathbb{Z} & \text{Klein'sche Vierergruppe } V_4 & \text{Restklassen-Dichte } (E,A,B,C) \\
			\text{Ebene 2} & \text{Lipschitz-Gitter } \mathbb{H} & \text{Hyperoktaedergruppe } B_4 & \text{Trägergewicht } s_{\min} \in \{2,3,4\} \\
			\text{Ebene 3/4} & \text{Imaginäre Oktonionen } \operatorname{Im}\mathbb{O} & \text{Exzeptionelle Lie-Gruppe } G_2 & \text{Frame-Energie } 3p^2 \\
			\text{Ebene 5} & \text{Albert-Jordan-Algebra } h_3(\mathbb{O}) & \text{Exzeptionelle Lie-Gruppe } E_6 & \text{Kubische Determinante } \det(X) \\
			\text{Ebene 6} & \text{Freudenthal-System } F(h_3(\mathbb{O})) & \text{Exzeptionelle Lie-Gruppe } E_7 & \text{Quartische Norm } q(x) \\
			\text{\textbf{Ebene 7}} & \text{\textbf{Ausnahmegitter }} E_8 & \text{\textbf{H"{o}chste Symmetrie }} E_8 & \text{\textbf{Epstein-Zeta-Pol }} \operatorname{Res}_{s=8} = \frac{\pi^4}{3} \\
		\end{array}
	\end{equation}
	
	Die Symmetrieleiter ist somit ein selbstkonsistenter Formalismus, bei dem das $E_8$-Gitter als die fundamentale Quelle dient, deren niedrigdimensionale Schatten über die Jordan- und Freudenthal-Mittelungen schrittweise eine mathematische Verwandtschaft mit den Eichstrukturen aufweisen (Ebene C, heuristisch).
	
	\subsection{Formaler Demarkations-Schild (Ebene [C])}
	Wir halten mit absoluter mathematischer Unerbittlichkeit fest:
	Die universelle Optimalität des $E_8$-Gitters nach Cohn--Kumar--Miller--Radchenko--Viazovska sowie die Faktorisierung der Epstein-Zeta-Funktion sind **rigorose, unumstößliche Sätze der reinen Mathematik und algebraischen Zahlentheorie (Ebene A)**. Die Interpretation des Pols bei $s=8$ als physikalische Quelle oder die Projektion der Lie-Gruppen auf physikalische Eichbosonen des realen Universums ist eine **rein strukturelle und heuristische Analogie (Ebene C)**. Sie impliziert **keine Modifikation der Einstein'schen Feldgleichungen**, keine realen physikalischen Extradimensionen im makroskopischen Raumzeit-Kontinuum und keine Vorhersage von realen außergewöhnlichen GUT-Partnerteilchen im realen Laborraum.
	
\end{document}